\documentclass[11pt]{article}
\usepackage[utf8]{inputenc}
\ifdefined\pdfglyphtounicode
\input glyphtounicode.tex
\pdfgentounicode=1
\pdfglyphtounicode{parenleftbig}{0028}
\pdfglyphtounicode{parenleftBig}{0028}
\pdfglyphtounicode{parenleftbigg}{0028}
\pdfglyphtounicode{parenleftBigg}{0028}
\pdfglyphtounicode{parenrightbig}{0029}
\pdfglyphtounicode{parenrightBig}{0029}
\pdfglyphtounicode{parenrightbigg}{0029}
\pdfglyphtounicode{parenrightBigg}{0029}
\pdfglyphtounicode{bracketleftbig}{005B}
\pdfglyphtounicode{bracketleftBig}{005B}
\pdfglyphtounicode{bracketleftbigg}{005B}
\pdfglyphtounicode{bracketleftBigg}{005B}
\pdfglyphtounicode{bracketrightbig}{005D}
\pdfglyphtounicode{bracketrightBig}{005D}
\pdfglyphtounicode{bracketrightbigg}{005D}
\pdfglyphtounicode{bracketrightBigg}{005D}
\pdfglyphtounicode{floorleftbig}{230A}
\pdfglyphtounicode{floorleftBig}{230A}
\pdfglyphtounicode{floorleftbigg}{230A}
\pdfglyphtounicode{floorleftBigg}{230A}
\pdfglyphtounicode{floorrightbig}{230B}
\pdfglyphtounicode{floorrightBig}{230B}
\pdfglyphtounicode{floorrightbigg}{230B}
\pdfglyphtounicode{floorrightBigg}{230B}
\pdfglyphtounicode{ceilingleftbig}{2308}
\pdfglyphtounicode{ceilingleftBig}{2308}
\pdfglyphtounicode{ceilingleftbigg}{2308}
\pdfglyphtounicode{ceilingleftBigg}{2308}
\pdfglyphtounicode{ceilingrightbig}{2309}
\pdfglyphtounicode{ceilingrightBig}{2309}
\pdfglyphtounicode{ceilingrightbigg}{2309}
\pdfglyphtounicode{ceilingrightBigg}{2309}
\pdfglyphtounicode{braceleftbig}{007B}
\pdfglyphtounicode{braceleftBig}{007B}
\pdfglyphtounicode{braceleftbigg}{007B}
\pdfglyphtounicode{braceleftBigg}{007B}
\pdfglyphtounicode{bracerightbig}{007D}
\pdfglyphtounicode{bracerightBig}{007D}
\pdfglyphtounicode{bracerightbigg}{007D}
\pdfglyphtounicode{bracerightBigg}{007D}
\pdfglyphtounicode{angbracketleftbig}{27E8}
\pdfglyphtounicode{angbracketleftBig}{27E8}
\pdfglyphtounicode{angbracketleftbigg}{27E8}
\pdfglyphtounicode{angbracketleftBigg}{27E8}
\pdfglyphtounicode{angbracketrightbig}{27E9}
\pdfglyphtounicode{angbracketrightBig}{27E9}
\pdfglyphtounicode{angbracketrightbigg}{27E9}
\pdfglyphtounicode{angbracketrightBigg}{27E9}
\pdfglyphtounicode{slashbig}{002F}
\pdfglyphtounicode{slashBig}{002F}
\pdfglyphtounicode{slashbigg}{002F}
\pdfglyphtounicode{slashBigg}{002F}
\pdfglyphtounicode{backslashbig}{005C}
\pdfglyphtounicode{backslashBig}{005C}
\pdfglyphtounicode{backslashbigg}{005C}
\pdfglyphtounicode{backslashBigg}{005C}
\pdfglyphtounicode{radicalbig}{221A}
\pdfglyphtounicode{radicalBig}{221A}
\pdfglyphtounicode{radicalbigg}{221A}
\pdfglyphtounicode{radicalBigg}{221A}
\pdfglyphtounicode{summationtext}{2211}
\pdfglyphtounicode{summationdisplay}{2211}
\pdfglyphtounicode{producttext}{220F}
\pdfglyphtounicode{productdisplay}{220F}
\pdfglyphtounicode{integraltext}{222B}
\pdfglyphtounicode{integraldisplay}{222B}
\pdfglyphtounicode{uniontext}{222A}
\pdfglyphtounicode{uniondisplay}{222A}
\pdfglyphtounicode{intersectiontext}{2229}
\pdfglyphtounicode{intersectiondisplay}{2229}
\pdfglyphtounicode{coproducttext}{2210}
\pdfglyphtounicode{coproductdisplay}{2210}
\pdfglyphtounicode{contintegraltext}{222E}
\pdfglyphtounicode{contintegraldisplay}{222E}
\pdfglyphtounicode{unionsqtext}{2294}
\pdfglyphtounicode{unionsqdisplay}{2294}
\pdfglyphtounicode{unionmultitext}{228E}
\pdfglyphtounicode{unionmultidisplay}{228E}
\pdfglyphtounicode{logicalandtext}{2227}
\pdfglyphtounicode{logicalanddisplay}{2227}
\pdfglyphtounicode{logicalortext}{2228}
\pdfglyphtounicode{logicalordisplay}{2228}
\pdfglyphtounicode{circledottext}{2299}
\pdfglyphtounicode{circledotdisplay}{2299}
\pdfglyphtounicode{circleplustext}{2295}
\pdfglyphtounicode{circleplusdisplay}{2295}
\pdfglyphtounicode{circlemultiplytext}{2297}
\pdfglyphtounicode{circlemultiplydisplay}{2297}
\fi
\usepackage{amsmath,amssymb,amsthm}
\usepackage[margin=1in]{geometry}
\usepackage{booktabs,longtable,array}
\usepackage{xcolor}
\usepackage{hyperref}
\hypersetup{colorlinks=true,linkcolor=blue!45!black,urlcolor=blue!45!black,
 citecolor=blue!45!black,
 pdftitle={Progress on Turyn's converse: general symmetry theorems and coverage below 10,000},
 pdfauthor={Jeffery Kline},
 pdfsubject={Pre-release draft; general theorems, exact coverage and novelty assessments}}
\newtheorem{theorem}{Theorem}[section]
\newtheorem{lemma}[theorem]{Lemma}
\newtheorem{proposition}[theorem]{Proposition}
\newtheorem{corollary}[theorem]{Corollary}
\newtheorem{conjecture}[theorem]{Conjecture}
\theoremstyle{definition}
\newtheorem{definition}[theorem]{Definition}
\theoremstyle{remark}
\newtheorem{remark}[theorem]{Remark}
\setlength{\emergencystretch}{2em}
\renewcommand{\labelitemi}{\ensuremath{\bullet}}
\title{Progress on Turyn's converse:\\
general symmetry theorems and coverage below 10,000}
\author{Jeffery Kline}
\date{21 October 2026\\\small Version 0.1.0 --- \href{https://doi.org/10.5281/zenodo.23545780}{doi:10.5281/zenodo.23545780}}
\begin{document}
\maketitle
% ASSEMBLED SECTION: opening
\begin{abstract}
A symmetric Turyn pair of odd length $t$ is a pair of symmetric $\pm1$
sequences, one with a zero at the origin, whose periodic autocorrelations
add to $q=2t-1$ at shift zero and to zero elsewhere. Turyn constructed such
pairs whenever $q$ is a prime power and asked whether the converse holds;
it remains open. We prove restrictions on the symmetries a pair can have and
combine them with classical multiplier theorems. The two principal results
are these. If a prime $p\ge7$ divides $t$ exactly once, no pair is fixed by
the decimations that act as squares on the $p$-component and trivially on
the rest. On any group $H\times C_3$ with $H$ abelian of odd order, no pair
is fixed by inverting the $C_3$ factor unless $H$ is trivial. Both apply
when $t$ is squarefree, where earlier conductor tests say nothing. Together with exact quotient tests, norm-lattice certificates, earlier
results, and a ray-class transport that excludes an infinite family $q=13p$,
they exclude 655 of the 663 non-prime-power $q<10000$ with
$q\equiv1\pmod4$ that are sums of two squares. Eight remain open, the smallest $3965$; a mod-16 identity reduces
that case from 198 to 132 unknown signs without settling it. For each
result we say what is new and what specializes known methods. Priority was
checked against a bounded set of sources only.
\end{abstract}

\tableofcontents

\section{Introduction}\label{sec:scope}
This paper follows \cite{Kline26}, which excluded many composite orders
below $2000$ and set up the quotient framework used here. That paper is not
amended. Everything in Sections~\ref{early-section}--\ref{sec:phase} is new
relative to it; Section~\ref{sec:novelty} compares each result with the older
literature. Section~\ref{sec:contents} below summarizes what is proved, what
is a finite certificate, and what is new.

The finite statement is the following. It rests on the union of the criteria
and certificates listed in Section~\ref{sec:coverage}, not on one theorem.
\begin{theorem}[Recorded coverage below 10000]\label{thm:coverage}
Let
\[
 \mathcal E=\{q:5\le q<10000,\ q\equiv1\pmod4,\ q\text{ is not a prime
 power and is a sum of two squares}\}.
\]
Then $|\mathcal E|=663$. The unconditional arguments and exact certificates
listed in Section~\ref{sec:coverage} exclude symmetric cyclic Turyn pairs at
655 members of $\mathcal E$. They do not decide the remaining eight:
\[
 3965,\quad4745,\quad5185,\quad5525,\quad6409,\quad7565,\quad7685,\quad8245.
\]
\end{theorem}
These eight are open relative to the arguments collected here, not
necessarily in the whole literature. The exclusion of $5185$ known so far
assumes GRH and is not counted. Theorem~\ref{thm:coverage} concerns this
specific two-circulant construction, not arbitrary Hadamard matrices.

Within all 2499 parameters $q=5,9,\ldots,9997$, there are 642 prime powers,
where Turyn's construction gives existence; 1194 fail the necessary
two-squares condition; 655 have the further exclusions above; and eight
remain unresolved. These categories are disjoint. The degenerate parameter
$q=1$ is outside the stated range.

\subsection{What this paper contains}\label{sec:contents}
Most results take one form: a symmetry is shown to be impossible for any
pair, and a classical multiplier theorem then forces that symmetry at
particular $q$. T4 and the $C_3$ theorem need only a prime that divides $t$
exactly once, so they apply to squarefree $t$, where the full-conductor
test of Section~\ref{sec:conductor} holds automatically and says nothing. Table~\ref{tab:contents} lists the results, their kind, and
how each relates to earlier work. ``New'' means we did not find the
statement in the sources we checked (Section~\ref{sec:novelty}); it is not a
claim of first publication.

\begingroup\small
\setlength{\tabcolsep}{4pt}
\begin{longtable}{>{\raggedright\arraybackslash}p{0.34\textwidth}>{\raggedright\arraybackslash}p{0.15\textwidth}>{\raggedright\arraybackslash}p{0.43\textwidth}}
\caption{Results, their kind, and their relation to earlier work.}
\label{tab:contents}\\
\hline
Result & Kind & Relation to earlier work\\
\hline\endfirsthead
\hline Result & Kind & Relation to earlier work\\ \hline\endhead
Common multiplier group $H_0$ (D1) & theorem & Specialization of
 Delsarte--Goethals--Seidel \cite{DGS71}; not new.\\
Kernel injectivity; full conductor $F(4t,q)=4t$; cofactor threshold &
 theorems & Refines the scalar field-descent bound; new as stated.\\
Odd-index bound O; all-index prime-power bound T2 & theorems & Classical
 averaging; the nonvanishing step of T2 is new.\\
Full-unit fibre bound T1 & theorem & Orbit-energy method of Tan
 \cite{Tan18}; new as stated.\\
\textbf{Index-two prime-fibre theorem T4} & theorem & \textbf{Principal
 result}; new.\\
Odd-index fibre bound T5 & theorem & Short corollary of standard projection
 counting.\\
\textbf{Inverting a $C_3$ factor on $H\times C_3$} & theorem &
 \textbf{Principal result}; new.\\
Quotient energy QE; doubling-cycle tests & exact criteria & Classical
 contraction method \cite{JP99,AGZ21,Gordon26}; the criteria as used here
 are specific to this problem.\\
Norm-lattice certificates; colored catalogs; ray-class transport & certificates;
 theorem & Classical lattice and class-field tools; the transport of a
 certified obstruction is new.\\
Mod-16 identity; forced-fibre reconstruction & theorems & New as stated;
 prior-art search incomplete.\\
Coverage below $10000$ (Theorem~\ref{thm:coverage}) & finite record & 644
 exclusions by the uniform criteria above, 11 by earlier or case-specific
 certificates. These count coverage, not historically new exclusions.\\
\hline
\end{longtable}
\endgroup

We do not claim: the converse; any infinite family of excluded $q$ other
than the ray-class family $q=13p$ and the classical three-power family
$t=3^{4n+1}$ (Corollary~\ref{cor:three}); that the eight listed orders are open in
the whole literature; or an unconditional result at $5185$, where only a
GRH-conditional exclusion is known. Proofs and finite checks were written
with AI assistance and checked by separate AI reviewers and separately
written programs. That is not expert peer review.

\subsection{Objects and conventions}
Use additive notation for an odd finite abelian group $G$, let $0$ be its
identity, let $\delta$ denote the group-ring identity, and put
$J_G=\sum_{g\in G}g$. For an integer $m$ coprime to the exponent of $G$,
$\sigma_m$ sends $g$ to $mg$. The involution $X\mapsto X^{(-1)}$ inverts
group elements. On Gaussian group rings the adjoint $X^*$ also conjugates
coefficients. A real row is \emph{symmetric} when $X^{(-1)}=X$.

\begin{definition}
A symmetric Turyn pair on $G$, $|G|=t$ odd, consists of $A,B\in\mathbb Z[G]$
with $A_0=0$, $A_g\in\{-1,1\}$ for $g\ne0$, every $B_g\in\{-1,1\}$,
$A^{(-1)}=A$, $B^{(-1)}=B$, and
\begin{equation}\label{eq:turyn}
 A^2+B^2=(2t-1)\delta=q\delta.
\end{equation}
The original converse concerns the cyclic group $G=C_t$.
\end{definition}
Turyn proved existence when $q$ is a prime power and suggested the converse
\cite{Turyn72}. Throughout, ``order $q$'' refers to this parameter;
the row length is $t=(q+1)/2$, the corresponding negacyclic conference
matrix has order $q+1$, and the associated Hadamard construction has order
$2(q+1)$. A nonexistence conclusion for the pair does not transfer to all
Hadamard matrices of that order.

Putting $Z=A+iB$, symmetry turns \eqref{eq:turyn} into $ZZ^*=q\delta$.
Every character therefore has $|\chi(Z)|^2=q$. Augmentation gives
\begin{equation}\label{eq:augmentation}
 X^2+Y^2=q,\qquad X=A(1)\text{ even},\quad Y=B(1)\text{ odd}.
\end{equation}
If $\pi:G\to Q$ has odd fibre size $f$, its projected rows satisfy the
same norm equation. The projected $A$ origin is even and has absolute
value at most $f-1$; all other projected entries of $A$, and all entries
of $B$, are odd and bounded by $f$. Symmetry improves the origin condition
to $\bar A_0\equiv f-1\pmod4$. Quotient solutions need not lift to full
sign rows.

\subsection{The coefficient doubling identity}
\begin{lemma}\label{lem:doubling}
Every pair on an odd finite abelian group of order $t$ satisfies
\begin{equation}\label{eq:baseline-doubling}
 B_g=(-1)^{(t-1)/2}B_0A_gA_{2g}\qquad(g\ne0).
\end{equation}
Conversely, for any symmetric $A$ with the prescribed alphabet, choosing
$B_0\in\{-1,1\}$ and defining $B$ by this rule gives
$A^2+B^2\equiv q\delta\pmod8$ coefficientwise.
\end{lemma}
\begin{proof}
Write $A=J_G-\delta-2a$, $B=J_G-2b$, where $a,b$ are binary supports and
$a_0=0$, and let $c=(t-1)/2$. Expansion of \eqref{eq:turyn} gives
\[
 a^2+a+b^2=c\delta+\ell J_G,\qquad \ell=a(1)+b(1)-c.
\]
In characteristic two, squaring doubles group elements. The coefficient
at $2g$ is $a_g+a_{2g}+b_g=c[g=0]+\ell$. At $g=0$ it gives
$b_0=c+\ell$. Subtracting the latter relation at a nonzero $g$ and
returning to signs proves \eqref{eq:baseline-doubling}.
For the converse negate the whole $B$ row if necessary to take $B_0=1$.
Both binary supports then have even size. Put $K=J_G-\delta$ and
$D=\sigma_2$. The sign rule reads
$b=(c\bmod2)K+a+D^{-1}a$ in $\mathbb F_2[G]$. The exact defect is
\[
 4\bigl[a^2+b^2+a+(c-|a|-|b|)J_G-c\delta\bigr].
\]
Its bracket vanishes modulo two by Frobenius, the support parity and the
displayed expression for $b$. This proves the converse.
\end{proof}
This exact consequence is used by several later results. Its appearance
in a new application is not itself a claim of historical novelty.

% ASSEMBLED SECTION: early-theorems
\section{Multiplier, averaging, and prime-fibre obstructions}
\label{early-section}

Throughout this section, $t\geq3$ is odd, $q=2t-1$, and $A,B\in
\mathbb Z[C_t]$ are a symmetric Turyn pair. The identity group element is
$\delta$. For real coefficients, $*$ inverts group elements; for Gaussian
coefficients it also conjugates coefficients. Thus
\[
 A_0=0,\qquad A_j\in\{-1,1\}\ (j\ne0),\qquad
 B_j\in\{-1,1\},\qquad A^*=A,\quad B^*=B,
 \qquad A^2+B^2=q\delta.
\]
Write $\sigma_u(\sum_j c_jg^j)=\sum_j c_jg^{uj}$ for decimation by a unit
$u$. Every multiplier used below fixes both rows coefficientwise. Signed
multipliers, which may negate one row, need not satisfy this hypothesis.

\subsection{The common multiplier input}
\label{early-multipliers}

For every divisor $m>1$ of $t$, let $A^{[m]},B^{[m]}$ be the images under
the quotient map $C_t\longrightarrow C_m$. Define
\begin{equation}
 \mathcal H(m)=\bigcap_{r\mid q\;\text{prime}}
       \langle r^{e_r},-1\rangle\ \leq (\mathbb Z/m\mathbb Z)^\times,
 \qquad
 e_r=\begin{cases}1,&r\equiv1\pmod4,\\2,&r\equiv3\pmod4.\end{cases}
 \label{early-common-group}
\end{equation}
Put $\mathcal H_0=\mathcal H(t)$. The image of $\mathcal H_0$ modulo $m$
is contained in $\mathcal H(m)$, possibly strictly, so the quotient tests
may use the larger locally computed group. The groups in the intersection
need not be cyclic when $m$ has several prime factors.

\begin{proposition}[Common coefficient multipliers; D1]
\label{early-d1}
Every $u\in\mathcal H(m)$ fixes $A^{[m]}$ and $B^{[m]}$ coefficientwise.
\end{proposition}

\begin{proof}
This specializes classical common-multiplier reasoning; see the
Delsarte--Goethals--Seidel multiplier theorems
\cite[Theorems~5.2--5.4]{DGS71}. Set $X=A^{[m]}+iB^{[m]}$. Projection
preserves symmetry and the correlation identity, so $XX^*=q\delta$, and
inversion without coefficient conjugation fixes $X$.

By the Chinese remainder theorem $u$ has a unique lift $\widetilde u$
modulo $4m$ with $\widetilde u\equiv1\pmod4$. Let $s$ fix $i$ and invert
the $m$th roots of unity. Restricting $\langle r,s\rangle$ to
automorphisms fixing $i$ replaces $r$ by $r^{e_r}$, so $u\in\mathcal H(m)$
if and only if $\widetilde u\in\bigcap_{r\mid q}\langle r,s\rangle$.

At a character of order $d\mid m$, let $\alpha$ be the value of $X$ in
$\mathbb Q(i,\zeta_d)$. Then $\alpha\overline\alpha=q$ and $s$ fixes
$\alpha$. The primes dividing $(\alpha)$ lie above primes $r\mid q$, which
are unramified since $\gcd(q,4m)=1$. In this abelian cyclotomic extension
the Frobenius automorphism $\sigma_r$ fixes each prime ideal above $r$, so
it preserves the $r$-primary part of $(\alpha)$; $s$ preserves it because
$s(\alpha)=\alpha$. Hence $(\sigma_{\widetilde u}\alpha)=(\alpha)$, and
$\sigma_{\widetilde u}(\alpha)/\alpha$ is an algebraic integer.

Consider
\[
 E=q^{-1}\sigma_u(X)X^*\in\mathbb Q(i)[C_m].
\]
Every character value of $E$, at every character order including the
trivial character, is the preceding algebraic integer. By Fourier
inversion $mE_j$ is an algebraic integer in $\mathbb Q(i)$, so
$mE_j\in\mathbb Z[i]$. By definition $qE_j\in\mathbb Z[i]$. As
$\gcd(m,q)=1$, Bezout's identity gives $E_j\in\mathbb Z[i]$. Moreover
$EE^*=\delta$, whose identity coefficient is $\sum_j|E_j|^2=1$. Hence
$E=\varepsilon g^a$ with $\varepsilon$ a Gaussian unit. Inversion without
conjugation fixes $E$, so $2a=0$ in $C_m$, and $a=0$ since $m$ is odd. The
augmentation of $E$ is $1$, so $\varepsilon=1$. Finally
$\sigma_u(X)=EX=X$.
\end{proof}

The proof uses every character order. Invariance at primitive characters
alone would not give coefficientwise invariance on a composite quotient.

\subsection{Odd-index averaging and the prime-power bound}
\label{early-averaging}

Set $X^{[m]}=A^{[m]}+iB^{[m]}$ and $f=t/m$. At each nonzero position the
coefficient of $X^{[m]}$ has odd real and imaginary parts, each of
absolute value at most $f$. At the origin the real part is even with
absolute value at most $f-1$, and the imaginary part is odd with absolute
value at most $f$. Every character of $C_m$ takes a value of modulus
$\sqrt q$ on $X^{[m]}$. Let $\omega(m)$ be the number of distinct prime
divisors of $m$.

\begin{theorem}[Odd-index averaging; O]
\label{early-o}
Suppose $H\leq(\mathbb Z/m\mathbb Z)^\times$ fixes $X^{[m]}$ and has odd
index $d$. Then
\begin{equation}
 m^2\leq 2^{2\omega(m)-1}d^2q.
 \label{early-o-bound}
\end{equation}
In particular one may take $H=\mathcal H(m)$ when its index in
$(\mathbb Z/m\mathbb Z)^*$ is odd.
\end{theorem}

\begin{proof}
Average the unit action:
\[
 Y=\frac1{\varphi(m)}\sum_{u\in(\mathbb Z/m\mathbb Z)^\times}
       \sigma_u(X^{[m]}).
\]
$H$ has $d$ orbits on the unit positions. If $S$ is the sum of one
coefficient from each orbit, then $Y_1=S/d$. Since $d$ is odd and each
summand has two odd coordinates, $|S|\geq\sqrt2$.

The value of $Y$ at a character of exact order $e\mid m$ depends only on
$e$; call it $Z_e$. It is an average of values of modulus $\sqrt q$, so
$|Z_e|\leq\sqrt q$. Fourier inversion at position $1$ and the Ramanujan
identity $c_e(1)=\mu(e)$ give
\[
 \frac{mS}{d}=\sum_{e\mid m}\mu(e)Z_e.
\]
At most $2^{\omega(m)}$ terms are nonzero. Thus
$m\sqrt2/d\leq2^{\omega(m)}\sqrt q$. No norm identity for $Y$ is needed.
\end{proof}

At even index the parity argument alone allows $S=0$. For a prime-power
quotient this can be ruled out.

\begin{theorem}[Prime-power quotient bound at every index; T2]
\label{early-t2}
Let $m=\ell^a\mid t$, where $\ell$ is an odd prime and $a\geq1$. Let $H$
fix $X^{[m]}$, and put $d=\varphi(m)/|H|$, $f_\ell=t/\ell$, and
\[
 c_d=\begin{cases}\sqrt2,&d\text{ odd},\\2,&d\text{ even},\end{cases}
 \qquad \beta_\ell=\sqrt{(f_\ell-1)^2+f_\ell^2}.
\]
Then
\begin{equation}
 c_d|H|\leq
 \min\left\{2\sqrt q\left(1-\frac1\ell\right),\;
                   \sqrt q+\beta_\ell\right\}.
 \label{early-primepower-bound}
\end{equation}
Equivalently, the first bound is $c_dm/d\leq2\sqrt q$. Hence
$m^2\leq d^2q$ for even $d$, and $m^2\leq2d^2q$ for odd $d$.
\end{theorem}

\begin{proof}
Let $S$ again be the sum over the $d$ unit-orbit coefficients, and write
$Z_1=X(1)$. Summing the unit coefficients gives
\begin{equation}
 |H|S=Z_1-X^{[\ell]}_0.
 \label{early-unit-sum}
\end{equation}
The omitted coefficients are exactly those at original indices divisible
by $\ell$, so the origin term is that of the $C_\ell$ quotient, not the
$C_{\ell^a}$ quotient.

The averaging identity of the preceding proof now has only two nonzero
Ramanujan terms:
\begin{equation}
 \frac{mS}{d}=Z_1-Z_\ell,
 \qquad
 Z_\ell=\frac1{\ell-1}\sum_{\substack{\chi\text{ on }C_m\\
                                    \operatorname{ord}\chi=\ell}}
                                    \chi(X^{[m]}).
 \label{early-primepower-chord}
\end{equation}
If $S=0$, the average on the right equals $Z_1$, which has modulus
$\sqrt q$. Every term in the average has the same modulus, so equality in
the triangle inequality forces each term to equal $Z_1$. These terms are
exactly the nontrivial character values of $X^{[\ell]}$, and its trivial
character value is also $Z_1$. Fourier inversion would give
$X^{[\ell]}=Z_1\delta$. This is impossible, because each coefficient of
$X^{[\ell]}$ at a nonzero position has two odd coordinates. Therefore
$S\ne0$.

The coordinates of $S$ have the parity of $d$, so $|S|\geq c_d$: at even
index $S$ is a nonzero element of $2\mathbb Z[i]$. The triangle inequality
applied to \eqref{early-primepower-chord} gives the first bound. Applied
to \eqref{early-unit-sum}, with $|X^{[\ell]}_0|\leq\beta_\ell$, it gives
the second.
\end{proof}

\begin{remark}[Sharpness and scope]
\label{early-t2-boundary}
At $t=5$, $q=9$, $m=5$, let $\chi_5$ be the quadratic-character element
and $J_5$ the all-ones element of $\mathbb Z[C_5]$. The identities
$\chi_5^2=5\delta-J_5$ and $(J_5-2\delta)^2=J_5+4\delta$ show that
$A=\chi_5$, $B=J_5-2\delta$ is a Turyn pair. The subgroup of squares in
$(\mathbb Z/5\mathbb Z)^*$ fixes both rows. It has order $2$ and index
$2$, and the second (``box'') bound $\sqrt q+\beta_\ell$ is attained:
$4=3+1$. The nonvanishing lemma needs a prime-power quotient: at $q=29$
there is a pair with $S=0$ on the composite quotient $m=15$. The
parameter $f_\ell=t/\ell$ in the box bound is essential.
\end{remark}

\subsection{Full-unit fibres and Ramanujan energy}
\label{early-full-fibre}

Suppose $t=hm$ with $\gcd(h,m)=1$ and $m>1$. In the decomposition
$C_t=C_h\times C_m$, the subgroup
\[
 \mathcal F_m=\{u\in(\mathbb Z/t\mathbb Z)^\times:u\equiv1\pmod h\}
\]
acts trivially on $C_h$ and by all units on $C_m$.

\begin{theorem}[Full-unit fibre bound; T1]
\label{early-t1}
If both rows are fixed by $\mathcal F_m$, and $w=\omega(m)$, then
\begin{equation}
 2hm-1\leq(2h-1)4^w+(2^w-1)^2.
 \label{early-full-fibre-bound}
\end{equation}
Thus a violation of this inequality, together with
$\mathcal F_m\leq\mathcal H_0$, excludes a pair.
\end{theorem}

\begin{proof}
For each divisor $d\mid m$, let $e_d=\sum_{\gcd(k,m)=d}g_m^k$. The unit
group is transitive on each such index set, and $e_m=\delta$. Therefore
\[
 A=\sum_{d\mid m}\alpha_d\otimes e_d,\qquad
 B=\sum_{d\mid m}\gamma_d\otimes e_d,
\]
where the coefficients of $\alpha_d,\gamma_d\in\mathbb Z[C_h]$ are signs
except for $(\alpha_m)_0=0$. At any primitive character $\psi$ of $C_m$,
$\psi(e_d)=\mu(m/d)$. Define
\[
 E=\sum_{d\mid m}\mu(m/d)\alpha_d,\qquad
 F=\sum_{d\mid m}\mu(m/d)\gamma_d.
\]
For every character $\phi$ of $C_h$, the pair equation at
$\phi\otimes\psi$ gives $|\phi(E)|^2+|\phi(F)|^2=q$. Parseval on $C_h$
then gives $\|E\|^2+\|F\|^2=q$.

Exactly $2^w$ terms survive in the defining sums, so every coefficient of
$E,F$ has absolute value at most $2^w$, except that $|E_0|\leq2^w-1$
because $A_0=0$. Summing the squared bounds proves
\eqref{early-full-fibre-bound}.
\end{proof}

For an odd coprime split, the inequality fails for every $m>1$ except
$m=3$, and $m=15$ with $h\geq7$ coprime to $15$. Indeed, for $w=1$ it
fails exactly when $m\geq5$; for $w=2$ it fails when $m\geq17$, and when
$m=15,h=1$. For $w\geq3$, the product of the first $w$ odd primes exceeds
$4^w$, which suffices. These are the cases the bound leaves open, not
existence claims. The split $h=3,m=15$ is inadmissible because it is not
coprime.

\begin{corollary}[Quotient version of T1]
\label{early-t1-quotient}
Let $t'\mid t$, $W=t/t'$, and $t'=h'm$ with $\gcd(h',m)=1$. If the
projected rows on $C_{t'}$ are invariant under
$\{u\bmod t':u\equiv1\pmod{h'}\}$, then
\[
 q\leq(2h'-1)4^{\omega(m)}W^2+
                     (2^{\omega(m)}W-1)^2.
\]
\end{corollary}
\begin{proof}
The projected coefficients have absolute value at most $W$, except that
the real origin coefficient has absolute value at most $W-1$. The proof
of Theorem~\ref{early-t1} applies with these bounds; in particular the
reduced real origin coefficient is at most $2^{\omega(m)}W-1$.
\end{proof}

For example, $q=4453=61\cdot73$ has $t=2227=131\cdot17$. Since
$73\equiv-61^{-1}\pmod{2227}$, $\mathcal H_0=\langle61,-1\rangle$. The
orders of $61$ modulo $17$ and $131$ are $16$ and $5$. The elements
$61^{5j}$ fix the $C_{131}$ coordinate and, as $\gcd(5,16)=1$, exhaust
the units on $C_{17}$. Thus the full $17$-fibre is forced, whereas
\eqref{early-full-fibre-bound} would require $4453\leq1045$.

\subsection{The index-two prime-fibre theorem}
\label{early-index-two}

Let $t=hp$ with $p$ an odd prime and $\gcd(h,p)=1$. Write $\mathcal F_p$
for the prime fibre, and $\mathcal F_p^{(2)}$ for its square subgroup:
it fixes $C_h$ and multiplies the $C_p$ coordinate by a quadratic residue.
On $C_p$, put
\[
 J=\sum_{k\in C_p}g_p^k,\qquad
 \chi=\sum_{k\ne0}\left(\frac{k}{p}\right)g_p^k.
\]

\begin{theorem}[Index-two prime-fibre obstruction; T4]
\label{early-t4}
If $p\geq7$, no symmetric Turyn pair of length $hp$ is fixed by
$\mathcal F_p^{(2)}$. Consequently, a pair is excluded whenever
\[
 [\mathcal F_p:\mathcal H_0\cap\mathcal F_p]\leq2
 \qquad(p\geq7,\ v_p(t)=1).
\]
There is no upper bound on $h$ in this statement.
\end{theorem}

\begin{proof}
Invariance makes each $C_p$ fibre constant on $\{0\}$, the nonzero
squares, and the nonsquares. Write these three values for $A$ as
$\alpha_j,\beta_j,\gamma_j$, and set
\[
 S_j=\frac{\beta_j+\gamma_j}{2},\qquad
 D_j=\frac{\beta_j-\gamma_j}{2},\qquad U_j=\alpha_j-S_j.
\]
Exactly one of $S_j,D_j$ is a sign and the other is zero. Define
$S',D',V$ for $B$ in the same way, with $V_j=\alpha'_j-S'_j$. Then
\begin{equation}
 A=U\otimes\delta+D\otimes\chi+S\otimes J,\qquad
 B=V\otimes\delta+D'\otimes\chi+S'\otimes J.
 \label{early-legendre-decomp}
\end{equation}
Here $\alpha_0=0$, every other $\alpha_j$ is a sign, and every
$\alpha'_j$ is a sign.

Let $\epsilon=(-1/p)$. Summing the quadratic character gives
\[
 J^2=pJ,\qquad \chi J=0,\qquad
 \chi^*=\epsilon\chi,\qquad \chi\chi^*=p\delta-J.
\]
Symmetry of the full rows gives $U^*=U$, $S^*=S$, $D^*=\epsilon D$, and
likewise for $B$. Since $\delta,\chi,J$ are linearly independent,
expanding the correlation identity and comparing their coefficients gives,
in $\mathbb Z[C_h]$,
\begin{align}
 UU^*+VV^*+p(DD^*+D'D'^*)&=q\delta,\label{early-legendre-energy}\\
 UD+VD'&=0.\label{early-legendre-cross}
\end{align}
The raw $\chi$-coefficient is $2(UD+VD')$; dividing by $2$ is valid in
this torsion-free group ring.

Let $D_A=\operatorname{supp}D$ and $D_B=\operatorname{supp}D'$, with
sizes $n_A,n_B$. Call the complementary positions type $S$, and let
$M=2h-n_A-n_B$ be the total number of type-$S$ positions in the two rows.
We bound $M$ from above and below.

For the upper bound, let $k_A$ count nonzero type-$S$ positions of $A$
where $\alpha_j=-S_j$, and let $k_B$ count type-$S$ positions of $B$
where $\alpha'_j=-S'_j$. Put $k=k_A+k_B$, and let $\eta=1$ if $0$ is type
$S$ for $A$, $\eta=0$ otherwise. The identity coefficient of
\eqref{early-legendre-energy} is
\[
 \|U\|^2+\|V\|^2+p(n_A+n_B)=2hp-1.
\]
At nonzero type-$D$ positions $U_j^2=1$; at nonzero type-$S$ positions it
is $0$ or $4$. At the origin it is $1$ if $\eta=1$ and $0$ otherwise.
Thus
\[
 \|U\|^2+\|V\|^2=n_A+n_B-1+2\eta+4k,
 \qquad (p+1)M=2h+2\eta+4k.
\]
Because $k\leq M-\eta$,
\begin{equation}
 (p-3)M\leq2h-2\eta\leq2h.
 \label{early-upper-M}
\end{equation}

For the lower bound, reduce \eqref{early-legendre-cross} in
$\mathbb F_2[C_h]$. The sign alphabet gives
$\overline U=\overline D+\delta$ and $\overline V=\overline{D'}$, so
\[
 (\overline D)^2+\overline D+(\overline{D'})^2=0.
\]
Squaring in this commutative ring of characteristic two doubles
positions. Since $h$ is odd, doubling is a bijection, and the equation
becomes the support identity
\begin{equation}
 D_A=2(D_A\mathbin\triangle D_B).
 \label{early-support-identity}
\end{equation}
Taking cardinalities gives $|D_A\cap D_B|=n_B/2$. Inclusion--exclusion
then gives $n_A+n_B/2\leq h$, so $n_A+n_B\leq3h/2$. As $h$ is odd,
\begin{equation}
 M\geq(h+1)/2.
 \label{early-lower-M}
\end{equation}
Combining \eqref{early-upper-M} and \eqref{early-lower-M} gives
$(p-7)h+(p-3)\leq0$, impossible for $p\geq7$.

For the arithmetic consequence, $\mathcal F_p$ is cyclic, so a subgroup
of index at most two contains its square subgroup. Proposition~\ref{early-d1}
then supplies the forbidden coefficient symmetry.
\end{proof}

The threshold is necessary: the pair at $p=t=5$, $q=9$ in
Remark~\ref{early-t2-boundary} is fixed by the square subgroup. In that
example the nonsquares negate $A$ and fix $B$, which shows why the
sign-$+1$ hypothesis is needed. Hence the exact row-fixing stabilizer of
any pair meets a prime fibre with $p\geq7$ in index at least three. This
constrains primes occurring only once in $t$, where the full-conductor
condition gives no restriction.

The energy estimate alone is not enough. Without the support identity,
\eqref{early-upper-M} gives only a bound that depends on the cofactor $h$.
The mod-$2$ equation forces a positive proportion of type-$S$ positions,
which makes the prime threshold independent of $h$.

\subsection{Odd-index prime fibres}
\label{early-odd-fibre}

For $k\mid p-1$, let $K_k\leq\mathcal F_p$ be the unique subgroup of
index $k$, and put $f=(p-1)/k$. Its nonzero orbits in a $C_p$ fibre are
the $k$ multiplicative cosets, each of size $f$.

\begin{theorem}[Odd-index prime-fibre bound; T5]
\label{early-t5}
Let $t=hp$ with $\gcd(h,p)=1$. If $k$ is odd and a symmetric Turyn pair
of length $t$ is fixed by $K_k$, then
\begin{equation}
 f\leq k+1,\qquad p\leq k^2+k+1.
 \label{early-odd-fibre-bound}
\end{equation}
Consequently a pair is excluded if
$k=[\mathcal F_p:\mathcal H_0\cap\mathcal F_p]$ is odd and
$p>k^2+k+1$.
\end{theorem}

\begin{proof}
Project the rows to $C_h$ by summing each $C_p$ fibre. The projected
coefficients are
\[
 \overline A_j=\alpha_j+fs_j,\qquad
 \overline B_j=\alpha'_j+fs'_j,
\]
where $s_j,s'_j$ are sums of $k$ signs. They are odd, so have absolute
value at least one. Thus $|\overline A_j|\geq f-1$ for $j\ne0$,
$|\overline A_0|\geq f$, and $|\overline B_j|\geq f-1$ everywhere.
The identity coefficient of the projected correlation equation gives
\[
 2h(kf+1)-1
 =\sum_j(\overline A_j^2+\overline B_j^2)
 \geq(2h-1)(f-1)^2+f^2.
\]
Expanding and dividing by $2f>0$ gives $h(f-k-2)\leq-1$. Integrality
yields $f\leq k+1$, which proves both claims. The arithmetic consequence
follows from Proposition~\ref{early-d1}.
\end{proof}

This gives $p\leq3$ at index one, $p\leq13$ at index three, and
$p\leq31$ at index five. The boundary $p=13$, $k=3$ is attained at
$h=1$, $q=25$: exact enumeration gives twelve symmetric pairs, all fixed
by $K_3=\{1,5,8,12\}$. At $f=k+1$ the slack in the displayed energy bound
is $2(h-1)(k+1)$, so attainment at $h=1$ does not settle $h>1$.

For example, at $q=4141=41\cdot101$ and $t=2071=19\cdot109$,
$\mathcal H_0\cap\mathcal F_{19}$ has order six and index three. Since
$19>13$, Theorem~\ref{early-t5} applies. This fibre index must not be
replaced by the index of the image of $\mathcal H_0$ modulo $19$: the
fibre subgroup also fixes the other factor, while the image need not.

\subsection{Relation to established methods}
\label{early-novelty-scope}

D1 is classical multiplier input. The orbit averaging in O and T2, the
Ramanujan reduction in T1, and the projection and second-moment argument
in T5 are established methods. The prime-power nonvanishing step and the
even-index constant $c_d=2$ extend the older odd-index estimate; they are
not a new multiplier theorem. In T4 the decisive combination is the
Legendre decomposition, the support identity
\eqref{early-support-identity}, and the opposing bounds on $M$. Related
decomposition, energy, and characteristic-two arguments occur separately
in the literature.

% ASSEMBLED SECTION: conductor
\section{Multiplier kernels and full conductor}\label{sec:conductor}
This section constrains the actual coefficient stabilizer, not an enlarged
local group, by character cancellation and Parseval followed by the
classical common-power multiplier theorem.

\begin{theorem}[No principal-unit kernel]\label{thm:kernel}
Let $t=hp^a$ be odd, with $p$ prime and $(h,p)=1$. For $1\le b<a$, a
symmetric Turyn pair on $C_t$ cannot be fixed by every unit satisfying
$u\equiv1\pmod h$ and $u\equiv1\pmod{p^b}$.
Consequently the exact common coefficient stabilizer $H$ injects under
reduction modulo $\operatorname{rad}(t)$, and
$|H|\mid\varphi(\operatorname{rad}(t))$.
\end{theorem}
\begin{proof}
Set $L=p^{a-b}$ and $N=hp^b$, so $t=NL$. The kernel consists of the units
$u=1+Nz$, $z\bmod L$. If $v_p(j)=s<a-b$, the orbit of $j\bmod t$ is
$j+hp^{b+s}\mathbb Z\bmod t$, of length $p^{a-b-s}$; indices divisible
by $L$ are fixed. A character of full $p$-primary order $p^a$ annihilates
every nonfixed orbit, so its value on a kernel-invariant row is its value
on the subsequence $C'_k=C_{Lk}$, $k\bmod N$.

The two subsequences have total coefficient energy $2N-1$. Each of the
$h\varphi(p^b)$ characters of $C_N$ of full $p$-primary order $p^b$
lifts to such a character, so its two squared magnitudes sum to $q$.
Parseval over \emph{all} characters of $C_N$ gives
\[
 h\varphi(p^b)q\le N(2N-1),\qquad (p-1)q\le p(2N-1).
\]
But $q=2NL-1\ge2Np-1$, and
$(p-1)(2Np-1)-p(2N-1)=2Np(p-2)+1>0$, a contradiction.

For the corollary, the full reduction kernel is a product of cyclic groups
of orders $p^{a-1}$ over the distinct primes dividing $t$. A nontrivial
intersection with $H$ contains an element of prime order $p$. Its other
prime-primary components are trivial, and its $p$ component generates the
unique order-$p$ principal-unit subgroup. This is the forbidden kernel
with $b=a-1$, so the intersection is trivial.
\end{proof}

Define, for $p^a\Vert t$,
\[
 B_p=\max_{r\mid q}\left\{v_p(r^{p-1}-1)
              +v_p(\operatorname{ord}_{4\operatorname{rad}(t)}(r))\right\},
 \qquad F(4t,q)=4\prod_{p^a\Vert t}p^{\min(a,B_p)}.
\]
This is the field-descent function in the present coprime odd setting,
with the conventions of Tan's Definition 2.10 \cite{Tan18}.
\begin{theorem}[Full-conductor necessity]\label{thm:fc}
A cyclic symmetric Turyn pair necessarily satisfies $F(4t,2t-1)=4t$.
\end{theorem}
\begin{proof}
Fix $p^a\Vert t$, let $h=t/p^a$, and put $d_r=\operatorname{ord}_{4h}(r)$.
Passing from a prime $\ell\ne p$ to $\ell^j$ multiplies a multiplicative
order only by an $\ell$-power, and the modulo-$p$ order divides $p-1$.
Hence $v_p(d_r)=v_p(\operatorname{ord}_{4\operatorname{rad}(t)}(r))$.
If $B_p<a$, odd-prime lifting shows that $r^{(p-1)d_r}$ generates, modulo
$p^a$, all principal units congruent to one modulo $p^{b_r}$, where
$b_r=v_p(r^{p-1}-1)+v_p(d_r)$. It is one modulo $4h$. Therefore every
residue congruent to one modulo $4h$ and $p^{B_p}$ is a power of every
prime divisor $r$ of $q$.

To apply the common-power theorem of Delsarte--Goethals--Seidel
\cite[Theorems~5.2--5.4]{DGS71}, work in
$\mathbb Z[x]/(x^{2t}+1)$ and put $y=-x^2$, $z=x^t$. Then $y^t=1$,
$z^2=-1$, and
\[
 c(x)=A(y)+zB(y),\qquad c(x)c(x^{-1})=q.
\]
The two terms have disjoint even and odd exponent supports, so $c$ is a
normalized negacyclic conference polynomial with one zero coefficient.
The cited theorem makes each common power $u$ a multiplier:
$c(x^u)=\varepsilon x^{k}c(x)$. The unique zero coefficient forces $k=0$.
Since $u\equiv1\pmod4$, $z^u=z$, so $A(y^u)=\varepsilon A$ and
$B(y^u)=\varepsilon B$; as $B(1)$ is odd, $\varepsilon=1$. Substitution
fixes $z$ and sends $y$ to $y^u$, so separating the parity supports
fixes both rows coefficientwise. Reduction modulo $t$ supplies the kernel
forbidden by Theorem~\ref{thm:kernel}, with $b=B_p\ge1$.
Thus $B_p\ge a$ for every $p$, proving the assertion. The 2-part of the
field-descent function is four: in Tan's definition the modulus is
divisible by $4$ and every prime $r\mid q$ is odd.
\end{proof}
This statement concerns the arithmetic function $F$, not the claim that
every individual character value has minimal conductor $4t$. For
squarefree $t$ it is automatic.

\begin{corollary}[Explicit three-primary threshold]\label{cor:three}
Write $t=h3^a$, $a\ge1$, $(h,3)=1$, and
$e(h)=\max(\{0\}\cup\{v_3(p-1):p\mid h\})$. If a pair exists and
$q=2t-1$ is composite, then
\begin{equation}\label{eq:three-threshold}
 2\cdot3^a\le h3^{2e(h)}.
\end{equation}
In particular, at $t=3^a$ and at $t=5\cdot3^a$, existence forces $q$ prime.
\end{corollary}
\begin{proof}
Put $\beta=\max_{r\mid q}v_3(r^2-1)$. The 3-part of the unit-group
exponent modulo $4h$ is at most $3^{e(h)}$, so full conductor gives
$a\le\beta+e(h)$. Choose a prime $r\mid q$ attaining $\beta$ and let
$s=q/r>1$. For $M=3^{\min(a,\beta)}$, since $q\equiv-1\pmod M$, the
factors $r,s$ lie in opposite classes $\pm1\pmod M$. Both are odd and
greater than one, even if not coprime, so their minimum sizes are $2M-1$
and $2M+1$. Thus $2h3^a=q+1\ge4M^2\ge4\cdot3^{2a-2e(h)}$, which is
\eqref{eq:three-threshold}. For $h=1,5$, $e(h)=0$ and the inequality
fails for every $a\ge1$.
\end{proof}
For example, $h=7,13$ have
$e(h)=1$ and force primality for $a\ge4$, while $h=11$ does so for
$a\ge2$. The class $t=3^{4n+1}$, $n\ge1$, is excluded unconditionally:
$q=2\cdot3^{4n+1}-1$ is divisible by five and exceeds five.

Older scalar field-descent bounds also give eventual primality for fixed
$h$. The two formulations differ: at $q=305$, $t=153$, $F(612,305)=204$
and $\varphi(204)=64$, so the scalar inequality $305\le204^2/64$ passes
although full conductor fails.

% ASSEMBLED SECTION: symmetry-energy
\section{Exact quotient energy and coupled doubling cycles}
\label{se-sec-quotients}

Projection to a quotient turns a sign pair into integer rows: the
identity coefficient gives an energy equation, and the other
coefficients give the full correlation equations.  Exact multiplier
orbits make both reductions finite and exact for their stated
relaxations.  The doubling identity \eqref{eq:baseline-doubling} then
couples the two rows.

Let $M\leq (\mathbb Z/t\mathbb Z)^\times$ contain $-1$ and fix both
rows.  Its orbits on $\mathbb Z/t\mathbb Z$ are
$O_0=\{0\},O_1,\ldots,O_{N-1}$.  Write $a_i,b_i$ for the row values on
$O_i$, so $a_0=0$ and every other variable is a sign.  The applications
below use the guaranteed common multiplier group $\mathcal H_0$.
For $m\mid t$, let $I_m$ be the image of $M$ modulo $m$, and let
$\mathcal C_m$ be its orbits on $\mathbb Z/m\mathbb Z$.  We call these
quotient classes.  Write $\pi_m$ for the induced group-ring projection.

\begin{theorem}[Exact quotient-energy criterion]
\label{se-thm-qe}
Every orbit $O_i$ projects onto one quotient class $c(i)$, and its fibre
size
\[
 N_i=\#\{x\in O_i:x\equiv j\pmod m\}=\frac{|O_i|}{|c(i)|}
 \qquad (j\in c(i))
\]
is independent of $j$.  Define the finite signed-sum sets
\[
 \mathcal S^A_c=
 \left\{\sum_{\substack{i:c(i)=c\\i\ne0}}N_i a_i:a_i\in\{-1,1\}\right\},
 \qquad
 \mathcal S^B_c=
 \left\{\sum_{i:c(i)=c}N_i b_i:b_i\in\{-1,1\}\right\}.
\]
If a symmetric Turyn pair exists, there are
$X_c\in\mathcal S^A_c$, $Y_c\in\mathcal S^B_c$ satisfying
\begin{equation}
 \sum_{c\in\mathcal C_m}|c|(X_c^2+Y_c^2)=q.
 \label{se-eq-qe}
\end{equation}
Conversely, every solution of \eqref{se-eq-qe} in these finite sets is
realized by symmetric $M$-invariant alphabet rows whose projected energy
is $q$.  Thus absence of such a solution excludes a pair.
\end{theorem}

\begin{proof}
The projection is equivariant for $M\to I_m$, so an orbit maps onto one
class.  If $j,j'$ lie in that class, some $u\in M$ takes $j$ to $j'$
modulo $m$, and multiplication by $u$ bijects the corresponding fibres
inside $O_i$.  Hence $\pi_m(A)$ and $\pi_m(B)$ have class values
\[
 X_c=\sum_{i:c(i)=c}N_i a_i,
 \qquad Y_c=\sum_{i:c(i)=c}N_i b_i.
\]
Projection preserves the group-ring equation, and its identity
coefficient is the sum of the squares of all projected coefficients;
this is \eqref{se-eq-qe}.  For the converse, each orbit lies in exactly
one class, so the signed sums in different classes, and in the two rows,
involve disjoint variables and can be chosen independently.  Since
$-1\in M$, the resulting rows are symmetric.  This converse asserts the
energy equation only.
\end{proof}

The complete projected identity is a stronger finite condition.  Put
\[
 M^{(r)}_{cc'}=\#\{j\in c:j+r\in c'\},\qquad r\in\mathbb Z/m\mathbb Z.
\]
Then every pair satisfies
\begin{equation}
 \sum_{c,c'\in\mathcal C_m}M^{(r)}_{cc'}
       (X_cX_{c'}+Y_cY_{c'})=q\,[r=0]
 \quad\text{for every }r.
 \label{se-eq-qf}
\end{equation}
This groups the correlation sum by the classes of $j$ and $j+r$.  We
call \eqref{se-eq-qe} the quotient-energy condition and
\eqref{se-eq-qf} the full quotient condition.  Their feasible points are
marginals; feasibility does not give a Turyn pair.

\begin{corollary}[Odd-multiplicity energy floor]
\label{se-cor-floor}
For each of $\mathcal S^A_c$ and $\mathcal S^B_c$, consider the multiset
of fibre weights $N_i$ of the orbits it sums.  Let $w$ be the largest weight, let $n$ be its multiplicity,
and let $R$ be the sum of all remaining weights.  Set
\[
 d=\begin{cases}\max(0,w-R),&n\text{ odd},\\0,&n\text{ even},\end{cases}
\]
and set $d=0$ for an empty multiset.  Denote the two resulting bounds by
$d_A(c),d_B(c)$.  A pair is impossible if
\[
 \sum_c |c|\bigl(d_A(c)^2+d_B(c)^2\bigr)>q.
\]
\end{corollary}

\begin{proof}
An odd number of signs has an odd, hence nonzero, sum.  So the terms of
weight $w$ have absolute sum at least $w$, and the remaining terms at
most $R$.  Apply the triangle inequality and Theorem~\ref{se-thm-qe}.
In the $A$-origin class, $O_0$ is excluded from the multiset because its
coefficient is zero.
\end{proof}

At $m=1$, the criterion keeps the exact row sums.  If the nonzero orbit
sizes are $s_i$ and $g=\gcd_i s_i$, then
\begin{equation}
 A(1)\equiv t-1\pmod{2g},\qquad
 B(1)\equiv t-1\pm1\pmod{2g},\qquad A(1)^2+B(1)^2=q.
 \label{se-eq-augmentation}
\end{equation}
Here $A(1),B(1)$ denote augmentation.  Counting the orbits of each size
sharpens these congruences.  At $q=8281$, the nonzero orbit sizes are
two $20$'s, two $50$'s, and twenty $200$'s, so
$B(1)\equiv\pm1\pmod{20}$.  The representations
$8281=0^2+91^2=84^2+35^2$, with the parity of the two row sums, require
$B(1)\in\{\pm91,\pm35\}$, a contradiction.
At $q=4685$, the orbit inventory gives
$A(1)\equiv\pm2\pmod{20}$ and
$B(1)\equiv\pm1,\pm3\pmod{20}$.  The only representations are
$4685=14^2+67^2=62^2+29^2$; the first fails the $A$ congruence and the
second the $B$ congruence.

The energy floor excludes $q=9809$ at $m=9$.
In every nonzero quotient class the weights are fifteen $36$'s, one $4$,
and one $1$; the origin has the same weights for $B$ and omits the $1$
for $A$.  Thus the energy is at least
\[
 32^2+31^2+8\cdot2\cdot31^2=17361>9809.
\]
The exact criterion also excludes $4705$ at $m=13$ and $4525$ at $m=31$;
their finite signed-sum data are in the coverage certificates.

\subsection{An exact run-count reduction}
\label{se-subsec-runs}

Normalize $B_0=1$ and put $\varepsilon=(-1)^{(t-1)/2}$.
Doubling permutes the nonzero $M$-orbits, because it is invertible and
commutes with $M$.  For a cycle $C$ of this permutation, write $L_C$ for
its length and $a_0,\ldots,a_{L_C-1}$ for its cyclic sequence of $A$
signs.  Let $k_C$ count its negative signs and let $2r_C$ count its
changes of sign, including the closing edge.

\begin{theorem}[Simultaneous quotient signatures]
\label{se-thm-runs}
Choose any collection of divisors $m\mid t$ such that doubling fixes
every class in $\mathcal C_m$.  Equivalently, $2\in I_m$, with $m=1$
interpreted trivially.  A doubling cycle lying above a class $c$ has a
constant fibre weight, denoted $N_{m,C}$.
Proposed simultaneous class coefficients $X_{m,c},Y_{m,c}$ are realized
by symmetric $M$-invariant alphabet rows satisfying the full equation
modulo $8$ (to which Lemma~\ref{lem:doubling} applies, since its proof uses
only that congruence) if and only if, for every nonzero doubling cycle, integers
$k_C,r_C$ can be chosen with
\begin{equation}
 \begin{cases}
 r_C=0,\quad k_C\in\{0,L_C\},&\text{or}\\
 1\le r_C\le\lfloor L_C/2\rfloor,
 \quad r_C\le k_C\le L_C-r_C,
 \end{cases}
 \label{se-eq-runbounds}
\end{equation}
and, simultaneously for every chosen quotient and class,
\begin{align}
 X_{m,c}&=\sum_{C\text{ above }c}N_{m,C}(L_C-2k_C),
 \label{se-eq-runs-a}\\
 Y_{m,c}&=[c=\{0\}]+\varepsilon
       \sum_{C\text{ above }c}N_{m,C}(L_C-4r_C).
 \label{se-eq-runs-b}
\end{align}
The same variables $k_C,r_C$ occur at all quotients.
\end{theorem}

\begin{proof}
The edge signs $a_i a_{i+1}$ have product $1$, so the number of
transitions is even.  The node and edge sums are
\[
 \sum_i a_i=L_C-2k_C,
 \qquad \sum_i a_i a_{i+1}=L_C-4r_C.
\]
A word without transitions is constant.  With $2r_C>0$ transitions,
there are $r_C$ nonempty runs of each sign, which gives
\eqref{se-eq-runbounds}.  Conversely, split $k_C$ into $r_C$ positive
run lengths and $L_C-k_C$ into another $r_C$, then alternate the runs.
This realizes every allowed pair $(k_C,r_C)$, the constant cases
included.

If doubling fixes a quotient class, all orbits in a cycle lie above that
class and have the same size, hence the same fibre weight.  The
coefficient identity $B_j=\varepsilon A_jA_{2j}$ then gives
\eqref{se-eq-runs-a}--\eqref{se-eq-runs-b}.
For sufficiency, choose the cyclic words independently on the disjoint
cycles, assign their signs to the $M$-orbits, and define $B$ by this
identity.  The rows have the prescribed projections and, by
Lemma~\ref{lem:doubling}, satisfy the full equation modulo $8$.
\end{proof}

Doubling always fixes the origin class, so for every $m\mid t$ its
possible $B$ coefficients are exactly
\begin{equation}
 1+\varepsilon\left[\frac{t}{m}-1-
       4\sum_{C\text{ above }0} w_Cr_C\right],
 \qquad 0\le r_C\le\lfloor L_C/2\rfloor,
 \label{se-eq-origin-set}
\end{equation}
where $w_C$ is the size of an orbit in $C$.  Since every allowed
transition count is attained by a cyclic word,
\eqref{se-eq-origin-set} is exact for the $M$-invariant mod-$8$
relaxation.  A pair must also satisfy $\pi_m(B)_0^2\le q$, a finite
obstruction from the origin coefficient alone.  Equations
\eqref{se-eq-runs-a}--\eqref{se-eq-runs-b} with \eqref{se-eq-qf} keep
the stronger coupling between different quotients.

\begin{corollary}[A fixed origin fibre]
\label{se-cor-fixed}
Let $h\mid t$, put $m=t/h$, and suppose the image of $M$ modulo $h$
contains $2$.  Then every pair satisfies
\[
 \pi_m(B)_0=B_0[1+\varepsilon(h-1)],
 \qquad [1+\varepsilon(h-1)]^2\le q.
\]
No coprimality assumption on $h$ and $m$ is required.
\end{corollary}

\begin{proof}
Choose $u\in M$ with $u\equiv2\pmod h$.  For $j$ divisible by $m$,
$(u-2)j$ is divisible by $t$, so $A_{2j}=A_{uj}=A_j$.
Hence each of the $h-1$ nonzero coefficients of $B$ in that fibre is
$\varepsilon B_0$, and the origin adds $B_0$.  Apply the projected
energy equation.
\end{proof}

For $q=6437$, $t=3219=3\cdot29\cdot37$ and
\[
 M=\langle41,-1\rangle\cap\langle157,-1\rangle
       \le (\mathbb Z/3219\mathbb Z)^\times
\]
has order $72$.  Above the origin of $C_3$, its three nonzero doubling
cycles have $(L_C,w_C)=(14,72),(2,18),(7,4)$, of total size $1072$, and
$\varepsilon=-1$.  The first edge contribution
$72(14-4r)=144(7-2r)$ is an odd multiple of $144$; the rest lies between
$-56$ and $64$.  So the total edge sum is at least $88$ or at most
$-80$, and
\[
 |\pi_3(B)_0|\ge81,\qquad 81^2=6561>6437,
\]
a contradiction.  The orbit inventory is an exact arithmetic
certificate.

For $q=7625$, $t=3813=3\cdot31\cdot41$ and
$M=\langle5,-1\rangle\cap\langle61,-1\rangle$.
The congruences
\[
 5^3\equiv125\equiv-61^{19}\pmod{3813},
 \qquad125\equiv2\pmod{123}
\]
show that the image of $M$ modulo $123$ contains $2$.
As $\varepsilon=1$, Corollary~\ref{se-cor-fixed} gives
\[
 |\pi_{31}(B)_0|=123,
 \qquad123^2=15129>7625,
\]
with no enumeration of full quotient marginals.

\section{A uniform obstruction to inversion of a three-element factor}
\label{se-sec-three}

The energy method for a full prime fibre has an exceptional boundary at
$p=3$.  The next theorem resolves it for every odd finite abelian
complementary group, using the complete coefficient equations rather
than only their identity coefficients.  Throughout this section, $\delta$ is the identity element of the group
ring, $J_H=\sum_{g\in H}g$, and $\sigma_a$ is the automorphism induced
by multiplication by $a$ when this is invertible on $H$.

\begin{lemma}[Quadratic symmetry transfer]
\label{se-lem-transfer}
Let $H$ be finite abelian, let $E,\Lambda\in\mathbb Z[H]$, and suppose
\[
 \Lambda(3\Lambda-2E)=q\delta,\qquad 3\nmid q.
\]
Every group automorphism fixing $E$ also fixes $\Lambda$.
\end{lemma}

\begin{proof}
For an automorphism $\sigma$ fixing $E$ and a character $\phi$ of $H$,
the algebraic integers $\phi(\Lambda)$ and $\phi(\sigma\Lambda)$ are
roots of $3z^2-2\phi(E)z-q=0$.  If they were distinct, their product
would be $-q/3$, which is not an algebraic integer.  So they agree for
every character, and character separation gives $\sigma\Lambda=\Lambda$.
\end{proof}

\begin{lemma}[Complementary-support rigidity]
\label{se-lem-support}
Let $H$ have odd order $h$, and suppose $W$ and $2W$ partition
$H\setminus\{0\}$.  Let $D\in\mathbb Z[H]$ have support $W$ and
coefficients $\pm1$ on $W$.  If
$E=\epsilon\delta+2D$, where $\epsilon\in\{-1,1\}$, satisfies
$\sigma_2(E)^2=E^2$, then
\[
 E+\sigma_2(E)=2\epsilon J_H,
 \qquad E(1)=\epsilon h,
 \qquad D=\epsilon\mathbf1_W.
\]
\end{lemma}

\begin{proof}
The support partition gives
$\sigma_2(D)-D=J_H-\delta+2Q$ for some $Q\in\mathbb Z[H]$.
For a nontrivial character $\phi$, equal squares give
$\phi(\sigma_2 E)=\pm\phi(E)$.  The plus sign would give
$0=-1+2\phi(Q)$, contradicting algebraic integrality.  So every
nontrivial character annihilates $E+\sigma_2 E$, whose coefficients are
then constant by Fourier inversion; the identity coefficient is
$2\epsilon$.  This is the first formula, and augmentation gives the
second.  At $g\in W$ the first formula gives $2D_g=2\epsilon$, the
third.  When $h=1$, all conclusions are immediate.
\end{proof}

\begin{theorem}[Full three-fibre obstruction]
\label{se-thm-three}
Let $H$ be a finite abelian group of odd order $h$, and put
$G=H\times C_3$.  Suppose $A,B\in\mathbb Z[G]$ are symmetric,
$A_0=0$, all other $A$ coefficients and every $B$ coefficient are signs,
and
\[
 A^2+B^2=(6h-1)\delta.
\]
If both rows are fixed by the automorphism $(g,y)\mapsto(g,y^{-1})$,
then $h=1$.  At $h=1$ there are precisely the four pairs
\[
 A=(0,s,s),\qquad B=(b,-b,-b),\qquad s,b\in\{-1,1\}.
\]
\end{theorem}

\begin{proof}
Write
\[
 A=U+S(y+y^2),\qquad B=V+T(y+y^2).
\]
The four slices are symmetric on $H$, since the two nonzero $C_3$ slices
are equal and the rows are symmetric on $G$.  The trivial and
nontrivial characters of $C_3$ give the full identities
\begin{align}
 E^2+4F^2&=q\delta,
 &E&=U-S,&2F&=V-T,\label{se-eq-slice1}\\
 (U+2S)^2+(V+2T)^2&=q\delta,
 &&q=6h-1.\label{se-eq-slice2}
\end{align}
Here $E=\epsilon\delta+2D$, with $\epsilon=-S_0$,
$D_0=0$, and all coefficients of $D,F$ in $\{0,\pm1\}$.
Let $Z$ be the nonzero zero set of $D$, and let $Z'$ be the full zero
set of $F$.  The identity coefficient of \eqref{se-eq-slice1}, and
\eqref{se-eq-slice1} divided by $4$, give
\begin{equation}
 |Z|+|Z'|=\frac{h-1}{2},\qquad
 \epsilon D+D^2+F^2=c\delta,
 \quad c=\frac{3h-1}{2}.
 \label{se-eq-divided}
\end{equation}
By symmetry, the identity coefficients are sums of squares.

In characteristic $2$, squaring doubles subscripts, so
\eqref{se-eq-divided} reduces to
\begin{equation}
 \mathbf1_Z+\mathbf1_{2Z}+\mathbf1_{2Z'}
        =J_H+(c\bmod2)\delta.
 \label{se-eq-cover}
\end{equation}
Each of the $h-1$ nonzero elements is covered a positive odd number of
times on the left, while the total nonzero coverage is at most
\[
 2|Z|+|Z'|\le2(|Z|+|Z'|)=h-1.
\]
Equality forces $Z'=\varnothing$ and makes $Z,2Z$ a disjoint
partition of the nonzero elements.  The origin coefficient in
\eqref{se-eq-cover} forces $c$ odd, so $h\equiv1\pmod4$.  Hence
\begin{equation}
 V=-T=F,\qquad W:=\operatorname{supp}D=2Z,
 \qquad2W=Z,\qquad4W=W.
 \label{se-eq-tiling}
\end{equation}
In particular, $F$ is a full sign row.

Next we determine enough signs of $D$ to obtain a multiplier.  Write
$h=4k+1$ and let $N=\{s\in W:\epsilon D_s=-1\}$.  For a sign row $X$
supported on $Y$,
\[
 X^2=\mathbf1_{2Y}+2P(X),
\]
where $P(X)$ sums $X_aX_b$ over unordered distinct pairs with a fixed
sum; modulo $2$ it is the unsigned pair count $p_Y$.
For $Y=H$, every sum has $(h-1)/2=2k$ unordered distinct pairs, since
doubling is bijective.  Reducing \eqref{se-eq-divided} modulo $4$
therefore gives
\begin{equation}
 \mathbf1_N+p_W=J_H+(k+1)\delta
       \quad\text{in }\mathbb F_2[H].
 \label{se-eq-signrule}
\end{equation}
For $s\in W$, no diagonal pair in $W^2$ has sum $s$, because
$s/2\notin W$.  So, with
$c_W(s)=\#\{(a,b)\in W^2:a+b=s\}$,
\[
 D_s=\begin{cases}
 \epsilon,&c_W(s)\equiv2\pmod4,\\
 -\epsilon,&c_W(s)\equiv0\pmod4.
 \end{cases}
\]
By \eqref{se-eq-tiling} the bijection $(a,b)\mapsto(4a,4b)$ preserves
$W^2$, so $c_W(4s)=c_W(s)$ and
\begin{equation}
 \sigma_4(E)=E.
 \label{se-eq-Efour}
\end{equation}

Put $R=\sum_{s\in Z}S_s s$ in $\mathbb Z[H]$ and
$\Lambda=2R-\epsilon\delta$.  Coordinatewise,
\[
 U=\frac{\Lambda+E}{2},\quad
 S=\frac{\Lambda-E}{2},\quad
 U+2S=\frac{3\Lambda-E}{2},\quad V+2T=-F.
\]
Subtracting \eqref{se-eq-slice1} from four times
\eqref{se-eq-slice2} gives
\[
 9\Lambda^2-6\Lambda E=3q\delta,
 \qquad\Lambda(3\Lambda-2E)=q\delta.
\]
Since $3\nmid q$, Lemma~\ref{se-lem-transfer} and
\eqref{se-eq-Efour} imply $\sigma_4(\Lambda)=\Lambda$ and
$\sigma_4(S)=S$.

The coefficient doubling identity holds on every odd finite abelian
group: its proof uses only the sign encoding, characteristic-$2$
Frobenius, and invertibility of doubling.  On $G$ at the nonzero
coordinate $(s,y)$ the sign factor is $(-1)^{(3h-1)/2}=-1$, so
\[
 F_s=F_0S_sS_{2s}\qquad(s\in H).
\]
Since $S_{4s}=S_s$, this gives $F_{2s}=F_s$.  Applying $\sigma_2$ to
\eqref{se-eq-slice1} then gives $\sigma_2(E)^2=E^2$, and
Lemma~\ref{se-lem-support} gives $E(1)=\epsilon h$.

Finally let $f=F(1)$ and $v=(U+2S)(1)$.  Here $f$ is odd because $F$
has $h$ sign coefficients.  Augmenting the two slice identities gives
\[
 6h-1=h^2+4f^2=v^2+f^2.
\]
Since $f^2\ge1$, $(h-1)(h-5)\le0$; with $h\equiv1\pmod4$ this leaves
$h=1$ or $h=5$.  At $h=5$ the first equality forces $f^2=1$ and the
second gives $v^2=28$, impossible.  At $h=1$ the slice equations give
exactly the four stated pairs, and direct substitution verifies them.
\end{proof}

\begin{corollary}[Coordinate-free form]
\label{se-cor-involution}
Let $G$ be finite abelian of odd order, and let $\tau$ be an involutive automorphism
of $G$ with exactly three elements satisfying $\tau(g)=-g$.
A symmetric Turyn pair fixed by $\tau$ can exist only when $|G|=3$.
\end{corollary}

\begin{proof}
Division by $2$ is well-defined on $G$, and
\[
 g=\frac{g+\tau(g)}2+\frac{g-\tau(g)}2
\]
splits $G$ as the direct product of the fixed and anti-fixed subgroups.
Their intersection is trivial, being killed by $2$.  The anti-fixed
subgroup has order $3$, so $\tau$ is inversion on a $C_3$ factor and
the identity on its complement.  Apply Theorem~\ref{se-thm-three}.
\end{proof}

\begin{corollary}[Full prime-fibre obstruction]
\label{se-cor-prime}
Let $p$ be an odd prime and let $H$ be a finite abelian group of odd order $h$.  A symmetric
Turyn pair on $H\times C_p$ fixed by every automorphism acting
identically on $H$ and as a unit on $C_p$ exists only at $p=3,h=1$.
\end{corollary}

\begin{proof}
The case $p=3$ is Theorem~\ref{se-thm-three}.  For $p\ge5$, unit
transitivity on $C_p\setminus\{0\}$ gives
\[
 A=U+S(J_{C_p}-\delta),\qquad
 B=V+T(J_{C_p}-\delta).
\]
In the projection to $H$, the $A$-origin coefficient has magnitude
$p-1$, and every other $A$ coefficient and every $B$ coefficient has
magnitude at least $p-2$.  So the energy is at least
\[
 (p-1)^2+(2h-1)(p-2)^2,
\]
which exceeds $q=2ph-1$ by $2(p-1)[h(p-4)+1]>0$, a contradiction.
\end{proof}

For cyclic $t=ph$ with $p\nmid h$, the Chinese remainder theorem
identifies the relevant automorphisms with the pure $p$-fibre in
$(\mathbb Z/t\mathbb Z)^\times$.  So a guaranteed common multiplier
group containing this fibre excludes the parameter unless $t=3$.
For $p=3$, the pure fibre has order $2$; when $t>3$, its intersection
with the common row stabilizer must be trivial.

Theorem~\ref{se-thm-three} contains the earlier support-tiling and
mod-$16$ sign constraints without their finite search boundary, and
supersedes a bounded full-three-fibre exclusion for $1<h\le1700$.
Because $H$ is an
arbitrary odd abelian group, no cyclicity or coprimality assumption on
$H$ in $H\times C_3$ is needed; cyclic lengths with $9\mid t$ are not of
this form and are not covered.  The $p\ge5$ part of
Corollary~\ref{se-cor-prime} is the elementary energy argument; the
additional content is the $p=3$ boundary.  These results assume extra
symmetry; they are not the unrestricted Turyn converse.

% ASSEMBLED SECTION: catalogs
\section{Exact norm lattices and arithmetic transport}\label{sec:catalogs}
When the preceding criteria leave a nontrivial character problem, exact
integral lattices give a finite certificate without uncertified class-group
decisions. Their factorwise invariants also give an infinite, arithmetically
specified family.

\subsection{The period trace form and small-norm rigidity}
Let $\ell$ be an odd prime, $H\le\mathbb F_\ell^*$ contain $-1$, put
$h=|H|$, $d=(\ell-1)/h$, $K=\mathbb Q(\zeta_\ell)^H$, and $L=K(i)$.
For a generator $g$ of $\mathbb F_\ell^*$ write
$\eta_j=\sum_{a\in g^jH}\zeta_\ell^a$, $0\le j<d$.
\begin{lemma}\label{lem:period-trace}
The $\eta_j$ form an integral basis of $K$, $\mathcal O_L=\mathcal O_K[i]$,
and
\begin{equation}\label{eq:period-trace}
 T=(\operatorname{Tr}_{K/\mathbb Q}(\eta_j\eta_k))
   =\ell I_d-hJ_d.
\end{equation}
For $\alpha=\sum c_j\eta_j\in\mathcal O_K\setminus\mathbb Z$,
$\operatorname{Tr}(\alpha^2)\ge\ell-h$.
\end{lemma}
\begin{proof}
$H$ permutes the integral basis $\zeta_\ell,\ldots,\zeta_\ell^{\ell-1}$, so
the orbit sums are a basis of the invariant module $\mathcal O_K$. In a
product of two coset sums, exactly $h$ ordered pairs sum to zero when the
cosets agree, and none otherwise; the full cyclotomic trace is $\ell-1$ at
exponent zero and $-1$ otherwise. Dividing by $h$ gives
\eqref{eq:period-trace}. Its determinant $\ell^{d-1}$ is coprime to the
discriminant $-4$ of $\mathbb Q(i)$, so the coprime-discriminant compositum
theorem gives $\mathcal O_L=\mathcal O_K[i]$.

The trace form is
\[
 Q(c)=h\sum_{j<k}(c_j-c_k)^2+\sum_jc_j^2.
\]
Since $\sum\eta_j=-1$, $\alpha$ is rational exactly when all $c_j$ are
equal. Otherwise, splitting at an integer between their minimum and
maximum gives at least $d-1$ unequal pairs, and $\sum c_j^2\ge1$, so
$Q(c)\ge h(d-1)+1=\ell-h$.
\end{proof}

\begin{proposition}[Two rigidity thresholds]\label{prop:smallnorm}
Suppose $\alpha,\beta\in\mathcal O_K$ and $\alpha^2+\beta^2=N$ for a
positive integer $N$.
If $dN<\ell-h$, both components are rational integers. The same holds
at equality if $d$ is odd.
If, in addition, $\alpha+i\beta\equiv1\pmod{2\mathcal O_L}$, then a
stronger sufficient condition is
\[
 (d\text{ odd or }N<\ell),\qquad dN<8(\ell-h)+d-4.
\]
\end{proposition}
\begin{proof}
Taking traces gives the strict first bound. At equality, a nonrational
component would take the entire trace and make the other zero, so its
square would be rational; an odd-degree field has no quadratic subfield.

For the second bound, write the components in the period basis: all real
coefficients are odd and all imaginary coefficients even. A nonconstant
odd vector costs at least $4h(d-1)+d$ in the trace form, and a nonconstant
even vector at least $4h(d-1)+4$, so two nonconstant components cost at
least $8(\ell-h)+d-4$. If exactly one is nonrational, its square is an
integer. Odd degree again rules this out. In even degree the only possible
quadratic subfield of the real cyclic field is $\mathbb Q(\sqrt\ell)$;
a nonrational integral element with rational square is an integral multiple
of $\sqrt\ell$, of square at least $\ell$. This contradicts $N<\ell$.
\end{proof}
The degree-or-norm clause is essential. Forced Gaussian ideal factors may
first be removed from a character value, provided the residual is proved
integral with the right parity. At $8865$, with conductor 31, $d=5$, $h=6$,
removing factors of norms 9 and 197 leaves norm 5, at the endpoint
$dN=\ell-h=25$. The character is Gaussian. After division by three, its odd
real coordinate has residues $\{\pm27,\pm29\}$ modulo 31, whereas an even
real augmentation has residues $\{\pm12,\pm16\}$; these sets are disjoint.

\subsection{A complete exact enumeration domain}
\begin{theorem}[Trace-minimal generators]\label{thm:trace-minimal}
Let $L$ be a CM field of degree $2d$, let $q$ be a positive rational integer,
and let $I$ be an integral ideal satisfying $I\overline I=(q)$. For every
nonzero $x\in I$,
\[
 \operatorname{Tr}_{L/\mathbb Q}(x\overline x)\ge2dq,
\]
with equality exactly when $x\overline x=q$.
\end{theorem}
\begin{proof}
The inclusion $(x)\subseteq I$ gives $|N_{L/\mathbb Q}(x)|\ge N(I)=q^d$.
Apply the arithmetic--geometric mean inequality to the $d$ positive
conjugates of $x\overline x$. Equality requires $|N(x)|=q^d$ and equal
conjugates, each then equal to $q$. The converse is immediate.
\end{proof}
In the period setting, a full integral basis matrix $B_I$ of $I$ gives the
normalized lattice form
\[
 Q_I=\frac1q B_I^{\mathsf T}\operatorname{diag}(T,T)B_I,
 \qquad \operatorname{rank}Q_I=2d,\quad
 \det Q_I=\ell^{2d-2}.
\]
The form is integral: $x\overline y/q\in\mathcal O_K[i]$ for $x,y\in I$,
so the trace of its real part is an integer. The determinant follows from
$|\det B_I|=N(I)=q^d$. By Theorem~\ref{thm:trace-minimal}, complete
enumeration through normalized length $d$ finds exactly the desired norm
generators. The determinant and cutoff do not depend on $q$, but the number
of allocation ideals can be large.

For $p^e\Vert q$, allocation on a pair $\mathfrak p,\overline{\mathfrak p}$
chooses exponents $a,e-a$, $0\le a\le e$; a self-conjugate prime requires
$2a=e$. Unique ideal factorization makes this exhaustive. An exported
lattice must be proved equal to the whole ideal by membership and exact
index; enumerating a proper sublattice would not certify nonexistence.

The $C_{61}$ computations use $H=\{1,13,14,47,48,60\}$, $h=6$, $d=10$. At
$8905$, all 16384 allocation ideals give 88 signed norm generators, 44
satisfying parity and eight boxed marginals with $(|X|,|Y|)=(72,61)$. The
complete $C_{73}$ classification has twelve marginals, all with
$(|X|,|Y|)=(16,93)$. A pair would have the same augmentation at both
quotients, a contradiction. So a nonempty marginal at one quotient does not
defeat the method: compatibility between quotients gives the exclusion.

At $2561=13\cdot197$, the $1048576$ allocations give 504 signed norm
generators, 252 satisfying parity; these have real conductor residues
\[
 \{0,\pm2,\pm8,\pm14,\pm16,\pm25,\pm30\}\pmod{61}.
\]
Every permitted augmentation has $(|X|,|Y|)=(40,31)$ or $(44,25)$, with
real residues $\{\pm21,\pm17\}$, disjoint from these. So no integral
augmentation-compatible lift exists. These classifications use exact
integer and rational arithmetic only; no unfinished class-group
certification is a premise.

\subsection{Factor catalogs and prescribed residue transformations}
\begin{theorem}[Colored factor matching]\label{thm:factor}
Let $I\overline I=(a)$, $J\overline J=(b)$ in a CM field, and suppose
$\gamma\overline\gamma=ab$, $(\gamma)=IJ$. Then multiplication by
$\gamma/b$ is an integral-lattice isometry from $\overline J$ to $I$,
with their respective normalized trace forms. If $a,b$ are odd,
$\gamma\equiv1\pmod2$, and a multiplicative residue map $\rho$ that
commutes with complex conjugation has $\rho(\gamma)=\alpha$ with
$b,\alpha$ invertible, their complete vector
counts agree under the color map
\[
 (x\bmod2,\rho(x))\longmapsto
 (\overline x\bmod2,\ (\alpha/b)\overline{\rho(x)})
\]
from $J$ to $I$.
\end{theorem}
\begin{proof}
As fractional ideals, $(\gamma/b)\overline J=IJ\overline J/(b)=I$, so the
map is bijective. The multiplier has relative norm $a/b$, which cancels the
normalization in the two forms. Its residue is one modulo two and
$\alpha/b$ under $\rho$. Composing with conjugation from $J$ to
$\overline J$ gives the stated color transformation.
\end{proof}
So a mismatch at an exactly enumerated length is an exclusion, while a
match on a finite prefix proves nothing about existence. For $2561$, each
factor has 56 Galois orbits of allocation ideals. Enumerating these 112
factor lattices through normalized length 80 rejects every product
allocation under all eight allowed augmentation transformations. An
independently organized exact enumeration confirms the mismatch and, with
the actual residues, keeps all 252 parity-compatible generators.

\subsection{Ray classes transport complete catalogs}
\begin{theorem}[Arithmetic transport]\label{thm:ray-transport}
Let $L/\mathbb Q$ be Galois and CM, of degree $2d$, with central complex
conjugation $c$. Let $\mathfrak m=N\mathcal O_L$, and let $p,p_0$ be
completely split rational primes coprime to $N$. Choose primes
$\mathfrak P,\mathfrak P_0$ above them with the same ray class modulo
$\mathfrak m$. Then the entire labeled collections of allocation lattices,
with normalized trace forms and ambient residue colors modulo
$\mathfrak m$, are isometric after a label permutation. If $i\in L$ and
$N$ is even, the Gaussian augmentation residues of norms $ap$ and $ap_0$
also agree for fixed $a$ coprime to $pp_0$.
\end{theorem}
\begin{proof}
Ray equality gives $\beta$ with $(\beta)=\mathfrak P/\mathfrak P_0$
and $\beta\equiv1\pmod{\mathfrak m}$ locally at the modulus primes.
For each set $T$ of one member from each pair $\{\sigma,c\sigma\}$
in $\operatorname{Gal}(L/\mathbb Q)$, set
\[
 I_p(T)=\prod_{\sigma\in T}\sigma\mathfrak P,
 \qquad\lambda_T=\prod_{\sigma\in T}\sigma\beta.
\]
Then $I_p(T)=\lambda_T I_{p_0}(T)$ and
\[
 \lambda_T\overline{\lambda_T}
 =\prod_{\sigma}\sigma\beta
 =N_{L/\mathbb Q}(\beta)=p/p_0.
\]
The last equality has no real-unit ambiguity: a CM field norm is positive,
and the fractional ideal norm is $p/p_0$. Multiplication by $\lambda_T$
is an integral-module bijection preserving the normalized trace form. It
preserves each ambient color, since $\lambda_T\equiv1$ at the modulus
primes and the difference of two integral vectors is integral at every
other prime. So the catalogs agree at \emph{every} length, not only up to
a cutoff.

For the Gaussian assertion, put $\Delta=N_{L/\mathbb Q(i)}(\beta)$.
If $\pi_0$ generates the prime below $\mathfrak P_0$, then
$\pi=\Delta\pi_0$ is a Gaussian integer of norm $p$, with
$\pi\equiv\pi_0\pmod N$. Unique factorization writes every norm-$ap$
Gaussian integer as $z\pi$ or $z\overline\pi$, where $N(z)=a$;
replacing $\pi$ by $\pi_0$ gives the residue-preserving bijection,
including the even-real, odd-imaginary parity requirement.
\end{proof}

\begin{corollary}[An infinite excluded family]\label{cor:ray-family}
Take
\[
 H=\{1,13,14,47,48,60\},\qquad K=\mathbb Q(\zeta_{61})^H,
 \qquad L=K(i),
\]
and put $\mathfrak m=122\mathcal O_L$. Fix a prime $\mathfrak P_0$ above 197.
Let $\mathcal S$ be the rational primes $p\ne13$, $p\nmid122$, which
split completely in $L$ and have some prime $\mathfrak P$ with
$[\mathfrak P]=[\mathfrak P_0]$ in $\operatorname{Cl}_{\mathfrak m}(L)$.
For every $p\in\mathcal S$ there is no cyclic symmetric Turyn pair at
$q=13p$, $t=(13p+1)/2$. The set $\mathcal S$ has positive natural density
among rational primes.
\end{corollary}
\begin{proof}
Taking norms of the ray congruence gives $p\equiv197\pmod{122}$, and
splitting gives $p\equiv1\pmod4$. So $61\mid t$, and the extended local
groups generated by $13,-1$ and by $p,-1$ modulo 61 both equal $H$.
The multiplier theorem puts the hypothetical primitive character in $L$,
with norm $13p$, parity one modulo two, and conductor residue equal to an
actual even-real, odd-imaginary augmentation.
Theorem~\ref{thm:ray-transport} carries all factor catalogs and augmentation
colors to those at $13\cdot197$, preserving the joining multiplier
$\alpha/p$ modulo 61. The certified empty joins through length 80 then
give a contradiction by Theorem~\ref{thm:factor}.

Let $M$ be the ray class field of $L$ for $\mathfrak m$. The modulus is
Galois-stable, so $M/\mathbb Q$ is Galois. Let $\mathcal C$ be the
conjugacy class of the Artin element of $\mathfrak P_0$ in
$\operatorname{Gal}(M/\mathbb Q)$. Its restriction to $L$ is trivial.
Chebotarev \cite[Theorems~28.6 and~28.9]{Sutherland21} identifies the completely split primes with the required ray
class precisely with this Frobenius class, up to finitely many omitted
primes. Hence their density is
\[
 \delta_{\mathcal S}=\frac{|\mathcal C|}{[M:\mathbb Q]}>0.
\]
\end{proof}
The class-field and Chebotarev inputs are classical; no GRH, numerical
class-group completeness or simultaneous-prime hypothesis is used. The
resulting $q$ are distinct squarefree composites with counting function
asymptotic to $\delta_{\mathcal S}\operatorname{Li}(X/13)$: positive
density among the varying prime factors, not among all integers. The
congruence $p\equiv197\pmod{244}$ alone is not claimed to characterize the
family.

A separately checked 77-digit prime in this ray state comes with an
explicit principal ray identity and ECPP/Pratt certificates. It certifies
membership only; it does not compute the ray class group or the density.
It lies beyond the range of Theorem~\ref{thm:coverage} and adds no row
below 10000.

% ASSEMBLED SECTION: phase
\section{A quadratic mod-16 criterion and forced reconstruction}\label{sec:phase}
The cycle counts of the preceding sections forget where the signs sit
within each cycle. The next congruence layer keeps those positions and
gives a uniform elimination theorem.

\subsection{The divided correlation defect}
Let $G$ be any odd finite abelian group of order $t>1$, normalize $B_0=1$,
and let $P$ be the negative support of $A$. Write $D=\sigma_2$,
$K_0=J_G-\delta$, $c=(t-1)/2$, $e=c\bmod2$ and
$\eta=|P|/2\bmod2$. For a binary support $X$, let $W(X)$ be the
group-ring sum of unordered pairs of distinct elements of $X$, reduced
modulo two. All products below are convolutions.

\begin{theorem}[Quadratic mod-16 identity]\label{thm:mod16}
Define $B$ by \eqref{eq:baseline-doubling}. Then the exact integer
correlation defect is divisible coefficientwise by eight, and in
$\mathbb F_2[G]$,
\begin{align}\label{eq:quadratic-F}
 F(P)&=W(P)+P D(P)+eP+(1-e)D(P)+D^2(P)\notag\\
 &\quad+\bigl(\lfloor c/2\rfloor+\eta\bigr)K_0+\eta\delta
 =D\left(\frac{A^2+B^2-q\delta}{8}\right).
\end{align}
Consequently the full mod-16 equations are exactly $F(P)=0$.
\end{theorem}
\begin{proof}
Put $Q=(1-B)/2$ coefficientwise, $R=D^{-1}P$,
$\epsilon=1-2e$, $p_0=|P|$ and $T_0=\sum_gP_gR_g$.
Both $p_0$ and $T_0$ are even by inversion symmetry. The sign rule gives
\[
 Q=eK_0+\epsilon(P+R-2P\circ R),\quad
 |Q|=e(t-1)+\epsilon(2p_0-2T_0),
\]
where $\circ$ is coefficientwise intersection. Integer expansion gives
\[
 A^2+B^2-q\delta=4[P^2+Q^2+P+\alpha J_G-c\delta],\quad
 \alpha=c-p_0-|Q|.
\]
For a binary vector $X$, $X^2=DX+2W_{\mathbb Z}(X)$, where
$W_{\mathbb Z}$ is the unreduced unordered-pair sum. At a nonzero $s$,
the remaining numerator $D(P+Q)_s+P_s+\alpha$ is even, with half modulo
two equal to
\[
 \lfloor c/2\rfloor+p_0/2+(1-e)(P_{s/2}+P_s)+P_{s/2}P_s.
\]
This uses $e+\epsilon c=2\epsilon\lfloor c/2\rfloor$ and even $T_0$. At
the origin the correlation defect is exactly zero.

For binary $X,Y$, expanding the symmetric difference gives
\[
 W(X\mathbin\triangle Y)=W(X)+W(Y)+XY+D(X\circ Y)
 \quad\hbox{in }\mathbb F_2[G].
\]
Apply this to $Q=P\mathbin\triangle R\mathbin\triangle eK_0$.
Symmetric supports have even weight, so $K_0X=X$ in characteristic two.
At nonzero positions $W(K_0)_s=c-1$, and $e(c-1)=0\pmod2$. Therefore
\[
 W(P)+W(Q)=D^{-1}W(P)+PR+D(P\circ R)+e(R+DP)
\]
away from the origin. The coefficientwise intersection term cancels
$P_{s/2}P_s$ in the previous half-numerator. Applying $D$ gives
\eqref{eq:quadratic-F} at nonzero positions. At zero, $W(P)_0=p_0/2$,
$(PDP)_0=T_0=0\pmod2$, and the added $\eta\delta$ cancels the first
term.
\end{proof}

Write $P=\sum_i x_iE_i$ over disjoint even-sized symmetry orbits. The
coefficient of $x_ix_j$, $i\ne j$, is $E_iE_j+E_iD(E_j)+E_jD(E_i)$, the
linear coefficient is
\[
 W(E_i)+E_iD(E_i)+eE_i+(1-e)D(E_i)+D^2(E_i)
       +( |E_i|/2\bmod2)J_G,
\]
and the constant is $(\lfloor c/2\rfloor\bmod2)K_0$. So the criterion is
quadratic in the orbit bits. Replacing $A$ by $-A$ leaves $B$ and the
defect unchanged.

\subsection{The finite-field reconstruction theorem}
Let $p\equiv5\pmod8$ be prime, $H\le\mathbb F_p^*$ have order
$h\equiv2\pmod4$, and $n=(p-1)/h$ be even. Assume that 2 generates
$\mathbb F_p^*/H$. Let $C\le H$ have odd order $h/2$. Consider symmetric
rows on $C_p\times C_3$ invariant under $C$ on the first coordinate.
Normalize $A_{(0,1)}=A_{(0,2)}=1$ and $B_{(0,0)}=1$ and use the
doubling rule. If $y$ generates $C_3$, then
\[
 P=U+Vy+V^{(-1)}y^2,\qquad U^{(-1)}=U,\quad U_0=V_0=0.
\]
The binary vector $U$ is $H$-invariant and $V$ is $C$-invariant. From
here on, a star on a field value denotes the involution induced by
inversion on $C_p$.

\begin{theorem}[Forced inner fibre]\label{thm:forced}
Choose a primitive $p$th root $\zeta$ in characteristic two. Put
$L_2=\mathbb F_{2^{2n}}$, $K_2=\mathbb F_{2^n}$ and
\[
 u=U(\zeta),\quad v=V(\zeta),\quad w=W(V)(\zeta),\quad v^*=v^{2^n}.
\]
The $C$ periods form a basis of $L_2$, and the $H$ periods form a self-dual
normal basis of $K_2$. If $v\notin K_2$, every solution of the full mod-16
equations has the unique possible inner value
\begin{equation}\label{eq:forced}
 u=vv^*+\frac{vw+v^*w^*}{v^3+(v^*)^3}.
\end{equation}
This reconstructs one binary $H$-invariant $U$; it need not satisfy the
remaining equations. Odd $|V|$ is a sufficient condition for $v\notin K_2$.
\end{theorem}
\begin{proof}
Write $g=h/2$ odd and
$\theta_i=\sum_{a\in2^iC}\zeta^a$, $0\le i<2n$.
Since 2 generates modulo $H$ and $n$ is even, it generates modulo $C$,
and $2^nC=-C$. Frobenius cycles the $2n$ periods, each of absolute trace
one. Counting zero and nonzero exponents in their trace products gives
\[
 \operatorname{Tr}_{L_2/\mathbb F_2}(\theta_i\theta_j)
       =1+[j=i+n],
\]
because the uniform nonzero multiplicity is $g-[j=i+n]$ and the zero
contribution is even. The Gram matrix is $J+R$, where $R$ exchanges paired
indices; since $(J+R)^2=I$ in characteristic two, the periods form a
basis. The same count for $\beta_i=\theta_i+\theta_{i+n}$ gives
$\operatorname{Tr}_{K_2/\mathbb F_2}(\beta_i\beta_j)=[i=j]$.
So evaluation is a bijection on the stated binary supports, even when
the cyclotomic polynomial is reducible over $\mathbb F_2$.

As $g$ is odd, $\operatorname{Tr}_{L_2/\mathbb F_2}(v)=|V|\pmod2$, and
elements of $K_2$ have zero trace from $L_2$; this proves the odd-weight
guard. Here $t=3p\equiv7\pmod8$, so the class-$y^2$ component of
\eqref{eq:quadratic-F}, evaluated at $\zeta$, is
\begin{equation}\label{eq:semilinear}
 v^*u^2+(v^*+v^2)u=w+v(v^*)^2+v^*+(v^*)^4.
\end{equation}
Its unordered-pair component is $W(V)+UV^{(-1)}$;
the three product components are $U D_xV$, $V^{(-1)}D_xU$ and
$VD_xV^{(-1)}$. The linear terms give the last two terms on the right.

Conjugate \eqref{eq:semilinear} and eliminate $u^2$. The determinant is
$d_0=v^3+(v^*)^3$. If it vanished, $v/v^*$ would be a cube root of unity.
All cube roots of unity lie in $K_2$ because $n$ is even, so conjugation would both fix and
invert the ratio, forcing it to be one, contrary to $v\notin K_2$. Thus
$d_0\ne0$. The remaining numerator expands to $vw+v^*w^*+vv^*d_0$, giving
\eqref{eq:forced}. The expression is fixed by the involution, so it lies
in $K_2$, and the period basis reconstructs exactly one possible $U$.
\end{proof}

\begin{proposition}[The complete residual test]\label{prop:residual}
In Theorem~\ref{thm:forced}, define
\[
 \alpha=1+v^2/v^*,\quad
 \beta=w/v^*+vv^*+1+(v^*)^3,\quad
 a_0=\alpha+\alpha^*,\quad b_0=\beta+\beta^*.
\]
Then $a_0\ne0$, $u=b_0/a_0$, and semilinear compatibility is equivalent to
\begin{equation}\label{eq:compatibility}
 b_0^2+a_0(\alpha^*\beta+\alpha\beta^*)=0\quad\hbox{in }K_2.
\end{equation}
The remaining primitive class-zero equation is
\begin{equation}\label{eq:classzero}
 W(U)(\zeta)+vv^*+d_0+u^3+u+u^4+1=0.
\end{equation}
These two tests and the projected mod-16 equation on $C_3$ are together
equivalent to the full mod-16 equation in this normalized sector.
\end{proposition}
\begin{proof}
Divide \eqref{eq:semilinear} by $v^*$ to obtain
$u^2+\alpha u=\beta$. Subtract its conjugate. Since
$a_0=d_0/(vv^*)\ne0$, this forces $u=b_0/a_0$.
Multiplication of the substituted residual by $a_0^2$ yields precisely
\eqref{eq:compatibility}. Expanding the class-zero component of
\eqref{eq:quadratic-F} gives \eqref{eq:classzero}; the constant is one
because $t\equiv7\pmod8$.

Each $C$-invariant component vanishing at $\zeta$ vanishes at every
nontrivial $p$th root, since $C$ and Frobenius together reach all
exponents. In characteristic two, a polynomial of degree at most $p-1$
with these zeros is $0$ or $J_{C_p}$; augmentation separates them, since
$p$ is odd. So the projected equation excludes $J_{C_p}$ in each
component. The other nonzero $C_3$ component is the conjugate of the
first, and Theorem~\ref{thm:mod16} completes the equivalence.
\end{proof}
Every prime $p\equiv5\pmod8$ admits $h=(p-1)/2$, $n=2$, because 2 is a
quadratic nonresidue, so the reconstruction criterion has infinitely many
parameter instances. This is not an infinite family of exclusions, and the
needed symmetry is not claimed at every associated order.

\subsection{The exact 3965 consequence and the next gap}
At $q=3965$, $t=1983=3\cdot661$, the common multiplier group has order ten.
Its 199 nonzero orbits form doubling cycles of lengths 132, 66 and one.
After fixing the last sign, the exact $C_3$ equations and congruences force
\[
 \bar A=(0,11,11),\qquad\bar B=(61,-1,-1),\qquad
 (k_{132},r_{132})=(65,33),\quad(k_{66},r_{66})=(33,18);
\]
the complete integer enumeration of these quotient possibilities is in the
certificates. So $|V|=325$, $|U|=330$ for every normalized candidate, and
Theorem~\ref{thm:forced} eliminates all 66 inner variables, reducing 198
bits to 132.

For the saved outer support, the forced inner support has weight 310 and
fails \eqref{eq:semilinear}. Independently, the integer correlations give a
132-by-66 affine system of coefficient rank 66 and augmented rank 67; a
listed XOR of 35 rows is $0=1$. So \emph{every} one of its $2^{66}$ inner
fillings is excluded. All 66 relative rotations of the saved two words also
fail, while keeping the exact $C_3$ data and full mod-8 equations. Other
outer words remain unclassified, so this does not close 3965 and does not
change Theorem~\ref{thm:coverage}.

A cheaper trace filter has a proved limitation. Put the 132-type
orbit bits $x_j$ in doubling order, pair indices $i,i+n$, and align the
inner bits $y_i$ by the projected $H$ periods. Set $a_i=x_i+x_{i+n}$,
$z_i=a_i+y_i$. Specialize the $C_3$ generator to one \emph{before}
evaluating and taking the absolute trace from $K_2$. Then
\begin{equation}\label{eq:trace-filter}
 \sum_i\bigl(x_ix_{i+n}+a_iy_i+z_iz_{i+1}\bigr)=0.
\end{equation}
Self-duality makes the cross-pair contribution
$\sum_i(x_ix_{i+n}+a_iy_i)$ and the $PDP$ contribution
$\sum_i z_iz_{i+1}$; the other trace terms cancel. For a fixed inner
weight $k_y$ and run count $r_y$, expansion makes this linear in $y$:
\[
 \sum_i(a_i+a_{i-1}+a_{i+1})y_i
 =k_y-r_y+\sum_i x_ix_{i+n}+\sum_i a_ia_{i+1}.
\]
If both word weights are odd and $n$ is even, the left sides summed over
all $n$ relative rotations equal one. Since $n$ is even, they cannot all
be equal, so the filter accepts some phases and rejects others and cannot
exclude an entire pair of word shapes. For the saved 3965 words it accepts
37 phases and rejects 29; all fail the full mod-16 equations.

Whether full mod-16 correlation forces even $|V|$ throughout this
normalized arithmetic class is open. At $p=101,h=10,n=10$, two separate
exhaustive implementations tested all 524288 odd outer supports: 620 pass
semilinear compatibility, and none passes \eqref{eq:classzero}. This is a
finite result only. A nonzero relative trace to $\mathbb F_4$ cannot
replace the full class-zero condition: an explicit compatible $p=101$
support has zero projected defect and zero relative trace but a nonzero
class-zero residual; its full integer correlations are in the
certificates.

The open problem is to control the weighted unordered-pair terms in the
two complete residual equations.

% ASSEMBLED SECTION: novelty
\section{Prior work and novelty assessment}
\label{sec:novelty}

A checked proof is not a claim of priority. We compared each result with
the sources listed below; that search was bounded. Throughout, ``new''
means that we did not find the statement in those sources. Priority is
unresolved for every result in this paper.

\subsection{Earlier work on the same problem}
A symmetric Turyn pair of length $t$ gives a negacyclic conference matrix of
order $2t=q+1$ and a cyclic relative difference set with parameters
\[
 (m,n,k,\lambda)=(2t,2,2t-1,t-1)
\]
in $C_{4t}$. Delsarte, Goethals and Seidel pose the prime-power conjecture
for negacyclic conference matrices in Section~6 of \cite{DGS71}.
Jungnickel and Pott treat the corresponding $(q+1,2,q,(q-1)/2)$ family,
the conjecture, and multiplier-based nonexistence searches
\cite[pp.~344--347]{JP99}. These are predecessors on the same question.
Field descent and weighing-matrix bounds \cite{Schmidt99,LS05,Tan18} are
broader methods that we specialize.

Multiplier theory, character inversion, quotient group rings,
intersection numbers \cite{Turyn65}, Gaussian periods, self-conjugacy and orbit
enumeration are established methods, and we claim none of them. The
quotient-energy criterion, for instance, is the classical contraction
method: projected coefficients are sums over fibres, multiplier orbits
constrain those sums, and the identity coefficient of the norm equation is
an exact sum of squares. Such equalities and orbit searches appear in
\cite{AGZ21,Gordon26}, and \cite{JP99} searched this family by multiplier
orbits. What is specific here is the two-row formulation, its coupling
constraints, and the certificates.

An excluded order need not be newly excluded. For example, $q=425$ is the
period-$852$ problem settled by Arasu and Voss \cite{AV99}. The coverage
counts in this paper measure what the collected arguments decide, not how
many orders are excluded for the first time.

\subsection{Result-by-result comparison}
Table~\ref{tab:novelty} gives, for each result, the closest earlier work we
found and our status label, using the same labels as Table~\ref{tab:contents}.

\begingroup
\small
\setlength{\tabcolsep}{4pt}
\renewcommand{\arraystretch}{1.13}
\begin{longtable}{>{\raggedright\arraybackslash}p{0.25\textwidth}>{\raggedright\arraybackslash}p{0.47\textwidth}>{\raggedright\arraybackslash}p{0.21\textwidth}}
\caption{Closest earlier work and status of each result.}
\label{tab:novelty}\\
\hline
Result & Closest earlier work found & Status \\
\hline
\endfirsthead
\hline
Result & Closest earlier work found & Status \\
\hline
\endhead
Common multiplier group $H_0$ (D1)
& Common-power multiplier theorems with their sign convention
  \cite[Theorems~5.2--5.4]{DGS71}.
& Not new: specialization. \\
Doubling identity (Lemma~\ref{lem:doubling})
& Elementary characteristic-$2$ computation.
& Not claimed as new. \\
Odd-index bound O; $3$-adic index formula
& Averaging, parity and Ramanujan-sum inversion are classical.
& New as stated. \\
Three-power and fixed-cofactor eventual primality
& Follow from the published scalar $F$-bound; a second route uses
  \cite{DGS71}.
& Not new: specialization. \\
Kernel injectivity; full conductor; sharp cofactor threshold
& Classical multipliers and Parseval. The scalar $F$-bound does not give
  all applications, and the Mann tests of \cite{AJMP95,JP99} do not imply
  it by the substitution in Section~\ref{sec:mann}.
& New as stated; Arasu--Ma comparison open. \\
Full-unit fibre bound T1; all-index prime-power bound T2
& Orbit reductions and multiplier-order bounds in the proofs of
  \cite[Theorems~4.6 and~4.10]{Tan18}.
& New as stated, including T2's nonvanishing step. \\
\textbf{Index-two prime-fibre theorem T4}
& Legendre decompositions with energy bounds, e.g.\ \cite[Theorem~4.3]{Tan18};
  none has the quadratic-residue stabilizer with the mod-$2$ support identity.
& \textbf{Principal result; new.} \\
Odd-index fibre bound T5
& Short projection and second-moment argument.
& Corollary of standard counting. \\
Auxiliary extensions (Appendix~\ref{app:auxiliary})
& Extend T2 and T4 under stated hypotheses.
& No separate comparison made. \\
Quotient energy QE; finite sign searches
& \cite{JP99,AGZ21,Gordon26}; a related Gaussian two-component form in
  \cite[Section~5]{Tan26}.
& Method classical; criteria specific. \\
\textbf{Inverting a $C_3$ factor on $H\times C_3$}
& Global-inversion theorems of \cite{AJP90} use a different action
  (Section~\ref{sec:inversion}); none gives the symmetry-transfer chain.
& \textbf{Principal result; new.} \\
Doubling-cycle and fixed-fibre tests
& Elementary consequences of the doubling identity.
& Method elementary; criteria specific. \\
Small-norm rigidity; trace-minimal enumeration
& Gaussian-period trace forms and ideal-lattice enumeration.
& Specific cutoffs; method classical. \\
Colored catalogs; ray-class transport
& Ideal lattices, ray classes, CM norm products and Chebotarev
  \cite{Sutherland21}.
& New as stated. \\
Mod-$16$ identity; forced-fibre reconstruction
& Gaussian periods and characteristic-$2$ elimination.
& New as stated; prior-art search incomplete. \\
\hline
\end{longtable}
\endgroup

Two infinite families of excluded $q$ appear. The three-power family
$t=3^{4n+1}$, $n\ge1$, follows from the published scalar $F$-bound; we
include it as a specialization. The ray-class theorem carries the certified
$13\cdot197$ obstruction to a set of primes $p$ of positive density among
primes, giving excluded $q=13p$; the set of such $q$ has density zero among
integers.

\subsection{Mann-type tests}\label{sec:mann}
Quotient the relative difference set by the odd subgroup of order $t/d$,
$d\mid t$. For a prime $r\mid q$, Result~1.8 of \cite{AJMP95} gives
\[
 r^f\equiv-1\pmod{4d}\ \Longrightarrow\ v_r(q)=2j\ \text{and}\ r^j\le t/d,
\]
and \cite[Result~3.10]{JP99} allows the trigger $\tau r^f\equiv-1\pmod{4d}$
for an available multiplier $\tau$. Its other branch is empty here because
$q^2-4t(t-1)=1$. These tests overlap the full-conductor criterion but do
not imply it by this substitution. At $q=305=5\cdot61$, $t=153$, the
multipliers guaranteed by \cite[Result~3.8]{JP99} are generated by $5$ and
$61$ modulo $612$ and are all $1\bmod4$, so neither trigger fires for any
$d\mid153$; full conductor excludes $305$ because $F(612,305)=204<612$.
This separates the displayed tests only; $305$ was excluded long before.

Proposition~2.2 of \cite{AJMP95} assumes a splitting relative difference
set, and the Turyn lift in $C_{4t}$ is nonsplitting, so it does not apply.
The 1994 erratum \cite{AJP94} replaces a scalar by an arbitrary group-ring
coefficient in the 1991 proof \cite{AJP91}, repairs parts~(a),(b), rejects
equation~(11) and doubts part~(c). The size bound above was reproved with
arbitrary coset coefficients. We have not matched every part of the 1991
theorem to the later statements.

\subsection{Global inversion and forbidden multipliers}\label{sec:inversion}
The inverse-multiplier theorems of \cite{AJP90} do not imply T4 or the
$C_3$-inversion theorem. In the lift to $C_4\times G$, row symmetry acts by
$(z,g)\mapsto(z,-g)$, whereas a global inverse multiplier acts by
$(z,g)\mapsto(z^{-1},-g)$. Partial inversion of a $C_3$ factor and the
square-fibre actions of T4 also fix $z$. Theorems~4.2 and~6.5 of
\cite{AJP90} in fact forbid a global inverse multiplier for a proper Turyn
lift, so applying them to row symmetry would wrongly exclude prime powers.
Its Theorem~5.3 uses character values, a Galois sign change, coset constancy
and coefficient rigidity, which is close to the last step of our $C_3$
proof; it does not supply the symmetry-transfer steps that precede it.
Miller's forbidden-multiplier theorem \cite{Miller25} concerns ordinary
difference sets and an order-$p$ multiplier on a $p^e$ component with
$e\ge2$, a different object and action.

\subsection{Sources we could not check}
The divisible-difference-set theorem of Arasu--Ma \cite{AM01} is the most
important unchecked comparison: the
\href{https://www.sciencedirect.com/science/article/pii/S0378375800002780}{publisher}
returned HTTP 403 and the
\href{https://corescholar.libraries.wright.edu/math/564/}{author repository}
holds only an abstract. The full 1991 Mann paper, some field-descent sources
and several textbooks were not fully examined, and the mod-16 reconstruction
has had only a preliminary prior-art search. These gaps limit every ``new'' label in Table~\ref{tab:novelty}. Priority
claims of \cite{Kline26} are inherited, not re-examined here.

% ASSEMBLED SECTION: coverage
\section{Coverage, certificates, and reproducibility}\label{sec:coverage}
The finite claims use the range $5\le q<10000$, $q\equiv1\pmod4$.
The eligibility restriction removes prime powers and parameters failing
the two-squares condition; it is not a density among all integers.
Table~\ref{tab:coverage} records the disjoint dispositions and the
specified selected union.
\begin{table}[ht]
\centering
\begin{tabular}{lrr}
\toprule
Disposition or comparison & $q<2000$ & $q<10000$\\
\midrule
All parameters $q=5,9,\ldots$ &499&2499\\
Prime powers: existence known &164&642\\
Excluded by the two-squares condition &211&1194\\
Eligible non-prime-power two-squares parameters &124&663\\
Excluded by the selected reusable-criterion union &122&644\\
Additional recorded exclusions outside that union &2&11\\
Total eligible parameters excluded &124&655\\
Unresolved unconditionally in this record &0&8\\
\bottomrule
\end{tabular}
\caption{Reconciled coverage. The selected-union and additional-exclusion
rows partition the eligible exclusions; neither is added again to the
total.}\label{tab:coverage}
\end{table}

The selected union has eleven members, named as in \texttt{data/coverage.csv}:
O (Theorem~\ref{early-o}), FC (Theorem~\ref{thm:fc}), T1
(Theorem~\ref{early-t1}, not its quotient corollary), T2
(Theorem~\ref{early-t2}), T4 (Theorem~\ref{early-t4} with D1), T5
(Theorem~\ref{early-t5}), QE (Theorem~\ref{se-thm-qe}), C2\_uniform
(Theorem~\ref{se-thm-three}), Full\_prime\_fibre
(Corollary~\ref{se-cor-prime}), Fixed\_fibre (Corollary~\ref{se-cor-fixed})
and Doubling\_origin (Theorem~\ref{se-thm-runs}). These overlap. In particular, their
individual counts cannot be added. The older bounded three-fibre
enumeration is superseded by the uniform theorem, rather than being
presented as its proof.

\begin{proof}[Proof of Theorem~\ref{thm:coverage}]
The exact ledger contains one row for every $q=5,9,\ldots,9997$.
Its factorizations determine the prime-power and two-squares predicates.
The verification script in \texttt{verify/} recomputes all eleven
criteria from the statements above at every parameter. It finds exactly
644 eligible exclusions, agrees with every flag in
\texttt{data/coverage.csv}, and excludes no prime power. The 19 eligible survivors of this selected union are
\[
\begin{gathered}
909,1845,2561,2665,3965,4745,5185,5525,6137,6409,\\
7461,7565,7685,7929,8245,8865,8905,9245,9265.
\end{gathered}
\]
Table~\ref{tab:additional} supplies recorded exclusions for eleven of
these. Removing those eleven leaves exactly the eight stated parameters.
The source census, baseline evidence and this join are separately pinned
in the accompanying materials. The accompanying check script
refactors every parameter and verifies the set difference and all displayed
counts; it does not re-prove the earlier certificates.
\end{proof}

\begin{table}[ht]
\centering\small
\begin{tabular}{rp{0.69\textwidth}l}
\toprule
$q$ & Additional recorded argument & Material\\
\midrule
909 & Full-unit $C_7$ Gaussian augmentation-ray test &M05--M06\\
1845 & Classical self-conjugacy obstruction &M06\\
2561 & Complete $C_{61}$ parity/augmentation census; separately colored factor matching &M07\\
2665 & Complete quotient-correlation spectrum &M06\\
6137 & Full-unit $C_{11}$ Gaussian augmentation-ray test &M05--M06\\
7461 & Full-unit $C_7$ Gaussian augmentation-ray test &M05--M06\\
7929 & Recorded degree-two period predicate at $C_{13}$ &M06\\
8865 & Forced factors, small-norm rigidity and disjoint augmentation residues at $C_{31}$ &M02\\
8905 & Incompatible complete $C_{61}$ and $C_{73}$ augmentation spectra &M02--M03, M07\\
9245 & Full-unit $C_{23}$ Gaussian augmentation-ray test &M05--M06\\
9265 & Complete residual norm and augmentation check at $C_{113}$ &M02\\
\bottomrule
\end{tabular}
\caption{The eleven exclusions outside the selected union. Several use
general criteria; this is not a list of eleven isolated new results.}
\label{tab:additional}
\end{table}

The baseline ledger distinguishes evaluated source predicates, exact
enumerations and source theorem references. Its review reconstructed the
lower-frontier cases and checked the broader bookkeeping, but did not
rerun every expensive inherited solver above that frontier. That scope
remains visible here, especially for the recorded $7929$ predicate.
The present coverage theorem relies on these credited source
determinations as well as the later criteria. Not every row was recomputed for this paper.

\subsection{Current unresolved parameters}
\begin{table}[ht]
\centering
\begin{tabular}{rrll}
\toprule
$q$ & $t=(q+1)/2$ & Factorization of $q$ & Status\\
\midrule
3965&1983&$5\cdot13\cdot61$&Open; stronger mod-16 reduction\\
4745&2373&$5\cdot13\cdot73$&Open\\
5185&2593&$5\cdot17\cdot61$&Open unconditionally\\
5525&2763&$5^2\cdot13\cdot17$&Open\\
6409&3205&$13\cdot17\cdot29$&Open\\
7565&3783&$5\cdot17\cdot89$&Open\\
7685&3843&$5\cdot29\cdot53$&Open\\
8245&4123&$5\cdot17\cdot97$&Open\\
\bottomrule
\end{tabular}
\caption{Parameters below $10000$ left open by the arguments collected here.
The earlier GRH-dependent result at 5185 remains separately recorded.}
\label{tab:open}
\end{table}
An order disappears from this table only when a complete argument applies
to every hypothetical pair at that order. A timeout, a violated equation
at one chosen row, or a failure of one relaxation witness is not sufficient.
This is why the mod-16 work at 3965 does not change the table.

\subsection{What the exact checks establish}
For lattice certificates, exact prime-kernel membership and determinants
identify the full ideal; exact inclusive-bound vector enumeration proves
completeness. For quotient certificates, every allowed row or signed sum
belongs to an explicitly finite domain. For the final coverage join,
independently checked arithmetic flags and earlier certificates
are combined by exact set operations. Corrupted inputs and known positive
pairs were used in the supporting audits to test more than successful
execution.

The file \texttt{data/coverage.csv} has one row per parameter, with its
factorization, eligibility, status, firing criteria, material identifiers
and the conditional note at 5185. The script \texttt{verify/verify\_coverage.py}
uses only the Python standard library and runs in about half a minute.
Certificates for the eleven exclusions of Table~\ref{tab:additional} are in
\texttt{certificates/}. The file \texttt{audit/README.md} lists the material
identifiers M00--M17 with SHA-256 hashes of the research files they denote;
those files are not distributed.

\appendix
\section{Retained auxiliary results and limits}\label{app:auxiliary}
The following results are proved in the materials (M12) but not used in
the main text. They are not additional rows of coverage unless the selected
census explicitly uses them.

\paragraph{General subgroup averaging (T3).}
For $H\le G\le(\mathbb Z/m\mathbb Z)^*$ with $[G:H]$ odd, let $G_e$
be its image modulo $e\mid m$, and put
$\eta_{e,C}=\sum_{v\in C}\zeta_e^v$. An $H$-invariant projected row
system with odd Gaussian coordinates at unit positions satisfies
\[
 \frac{m\sqrt2}{[G:H]}\le\sqrt q
 \sum_{e\mid m}\sum_{C\in(\mathbb Z/e\mathbb Z)^*/G_e}|\eta_{e,C}|.
\]
Average over $G$, use the odd orbit-sum lower bound at position one,
and group Fourier inversion by these cosets. This contains the full-unit
averaging argument and variants using only the odd-order part of a quotient
unit group. The exact scope and correction to early firing claims are in
the retained proof/audit pair.

\paragraph{Squarefree Jacobi-kernel extension (I2m).}
For squarefree $m\mid t$, $(m,h)=1$, $h=t/m$, and $w=\omega(m)$,
invariance under the kernel of the Jacobi character on the full $m$-fibre
necessarily gives
\[
 (h+1)(m-2^{w+1}+1)\le4h(2^w-1)^2+2.
\]
The character component supplies a support identity analogous to T4, and
the identity component supplies the energy bound. Violation is sufficient
for exclusion; the statement is not an equivalence between this inequality
and pair existence. The full proof, including its large-$w$ induction, is
retained in M12.

\paragraph{A large-prime index-four theorem (Q4).}
In the coprime split $t=hp$, with $h\ge3$ odd, $p$ prime,
$4\mid p-1$ and $f=(p-1)/4$, coefficientwise invariance under the
index-four subgroup is impossible if
\[
 p\ge2h+1,\qquad f(f-1)>2h,\qquad
 2h<(f-1)(h-1).
\]
For $h\ge11$, the first inequality implies the other two. The source
proof uses the exact quotient identity, sparse exceptional supports and
two applications of Frobenius parity. Its corrected primitive congruence
must be combined with the quotient congruence modulo eight; the primitive
equation alone is insufficient. This theorem supplies no extra current
frontier closure in the recorded census.

\paragraph{Limits and refuted shortcuts.}
The source record retains the following negative findings because they
prevent incorrect generalizations:
\begin{itemize}
\item The isolated quartic emptiness proposal has an explicit compatible
Gaussian witness at 93721; it cannot support a nonexistence claim.
\item Universal descent of all $C_{61}$ norm values to a quadratic
subfield is false: 480 of the 504 norm-2561 values lie outside that
quadratic subfield.
\item Extending the three-primary cofactor threshold naively to every
prime is false; the known pair at $t=25,q=49$ supplies a counterexample.
\item A finite-prefix match of factor catalogs, a feasible quotient or a
mod-16 solution is a necessary-condition witness, not an exact pair.
\item At 3965 the specified degree-two certificate family has 803 rows
in 19901 squarefree monomial columns, rank 795, and no constant one in
its span. This does not establish nonlinear feasibility. The quadratic
parts of its 199 original equations have rank 199.
\item The proposed relative-$\mathbb F_4$ trace obstruction is refuted
at $p=101$. The stronger odd-support conjecture remains open.
\end{itemize}
The precise next generalization target is the compatibility of the forced
inner fibre with both residual equations, not a larger search over an
already feasible weaker relaxation. The explicit current formulations,
counterexamples and unresolved weighted-carry dependency are indexed in M17.

\section{Evidence map}\label{app:evidence}
The identifiers below denote the research files behind each result. They
are not distributed, except for the certificates in \texttt{certificates/};
\texttt{audit/README.md} gives their SHA-256 hashes.
\begin{longtable}{lp{0.39\textwidth}p{0.43\textwidth}}
\toprule
ID & Subject & Evidence and scope\\
\midrule\endhead
M00--M01&Published baseline; foundations&Original source and initial mathematical dependencies; later records supersede earlier status labels.\\
M02--M03&Small norms; 8865, 9265, 8905&Proofs, residual classifications and the complete compatible-augmentation contradiction.\\
M04--M06&Local arithmetic; prime powers; ledger&Proofs, exact quotient checks, baseline bookkeeping and explicit inherited-source limits.\\
M07&Reusable $C_{61}$ and factor matching&General lattice cutoff; complete 2561 certificate; separate colored-factor verification.\\
M08&Arithmetic catalog transport&Universal proof, class-field sources and one independently checked ray-member certificate.\\
M09--M10&Averaging and full conductor&Proofs, exact comparisons, kernel theorem and cofactor threshold.\\
M11&Mann and subsequent sources&Erratum interpretation and direct prior-theorem substitutions, with access limits.\\
M12--M13&Prime-fibre generalizations&Corrected source proofs, audits and later independent arithmetic reconciliation.\\
M14--M15&Quotient energy and antecedents&Six finite closures, corrected counting formulas and superseded bounded three-fibre work.\\
M16&Abelian three-fibre theorem; final census&Full proof and audit; 6437/7625 certificates; exact final criterion union and eight-case frontier.\\
M17&Mod-16 reconstruction and latest source&Full proof/audit; 3965 fixed-outer certificate; independent $p=101$ census; AJP1990 comparison.\\
\bottomrule
\end{longtable}

% ASSEMBLED SECTION: bibliography
\begin{thebibliography}{99}

\bibitem{Turyn65}
R.~J. Turyn, Character sums and difference sets,
\emph{Pacific J. Math.} \textbf{15} (1965), 319--346.

\bibitem{Turyn72}
R.~J. Turyn, An infinite class of Williamson matrices,
\emph{J. Combin. Theory Ser.~A} \textbf{12} (1972), 319--321.

\bibitem{Kline26}
J.~Kline, Toward a converse of Turyn's theorem: obstructions to Turyn pairs
at composite orders, version 0.2.0 (2026),
\href{https://doi.org/10.5281/zenodo.21776269}{doi:10.5281/zenodo.21776269}.

\bibitem{DGS71}
P.~Delsarte, J.~M. Goethals and J.~J. Seidel,
Orthogonal matrices with zero diagonal.~II,
\emph{Canad. J. Math.} \textbf{23} (1971), 816--832.
\url{https://doi.org/10.4153/CJM-1971-091-x}.

\bibitem{AJP90}
K.~T. Arasu, D.~Jungnickel and A.~Pott,
Divisible difference sets with multiplier $-1$,
\emph{J. Algebra} \textbf{133} (1990), 35--62.
\url{https://www.sciencedirect.com/science/article/pii/002186939090067X}.

\bibitem{AJP91}
K.~T. Arasu, D.~Jungnickel and A.~Pott,
The Mann test for divisible difference sets,
\emph{Graphs Combin.} \textbf{7} (1991), 209--217.
\url{https://doi.org/10.1007/BF01787628}.
Full theorem and proof not obtained in the bounded comparison.

\bibitem{AJP94}
K.~T. Arasu, D.~Jungnickel and A.~Pott,
Erratum: The Mann test for divisible difference sets,
\emph{Graphs Combin.} \textbf{10} (1994), 207.
\url{https://doi.org/10.1007/BF02986666}.

\bibitem{AJMP95}
K.~T. Arasu, D.~Jungnickel, S.~L. Ma and A.~Pott,
Relative difference sets with $n=2$,
\emph{Discrete Math.} \textbf{147} (1995), 1--17.
\url{https://doi.org/10.1016/0012-365X(94)00226-9}.

\bibitem{JP99}
D.~Jungnickel and A.~Pott, Perfect and almost perfect sequences,
\emph{Discrete Appl. Math.} \textbf{95} (1999), 331--359.
\url{https://doi.org/10.1016/S0166-218X(99)00085-2}.

\bibitem{AV99}
K.~T. Arasu and N.~J. Voss,
Answering a question of Pott on almost perfect sequences,
\emph{Des. Codes Cryptogr.} \textbf{18} (1999), 7--10.
\url{https://doi.org/10.1023/A:1008316615026}.
The prior-result identification uses the inspected abstract and bibliographic
record; the complete proof was not reaudited.

\bibitem{AM01}
K.~T. Arasu and S.~L. Ma,
A nonexistence result on difference sets, partial difference sets and
divisible difference sets,
\emph{J. Statist. Plann. Inference} \textbf{95} (2001), 67--73.
\url{https://www.sciencedirect.com/science/article/pii/S0378375800002780}.
Metadata and abstract inspected; full DDS theorem unavailable in this comparison.

\bibitem{Schmidt99}
B.~Schmidt, Cyclotomic integers and finite geometry,
\emph{J. Amer. Math. Soc.} \textbf{12} (1999), 929--952.

\bibitem{LS05}
K.~H. Leung and B.~Schmidt, The field descent method,
\emph{Des. Codes Cryptogr.} \textbf{36} (2005), 171--188.
Inspected author version:
\url{https://personal.ntu.edu.sg/bernhard/Publications/pub/field.pdf}.

\bibitem{Tan18}
M.~M. Tan, Group invariant weighing matrices,
arXiv:1610.01914v4 (2018).
\url{https://arxiv.org/abs/1610.01914v4}.
The cited definition and theorem numbers refer to this inspected revision.

\bibitem{AGZ21}
K.~T. Arasu, D.~M. Gordon and Y.~Zhang,
New nonexistence results on circulant weighing matrices,
\emph{Cryptogr. Commun.} \textbf{13} (2021), 775--789.
\url{https://doi.org/10.1007/s12095-021-00492-0}.
Inspected preprint: \url{https://arxiv.org/abs/1908.08447v3}.

\bibitem{Gordon26}
D.~M. Gordon, Cyclic relative difference sets and circulant weighing matrices,
\emph{J. Algebraic Combin.} \textbf{63} (2026), Article~24.
\url{https://doi.org/10.1007/s10801-026-01515-w}.
Inspected preprint: \url{https://arxiv.org/abs/2501.14924v2}.

\bibitem{Miller25}
N.~Miller, Forbidden multipliers in abelian difference sets,
arXiv:2511.10231v1 (2025).
\url{https://arxiv.org/abs/2511.10231v1}.

\bibitem{Tan26}
M.~M. Tan, Character and multiplier obstructions for circulant weighing matrices,
arXiv:2608.18468v2 (2026).
\url{https://arxiv.org/abs/2608.18468v2}.
Used as a recent mechanism comparison, not as an independently audited
dependency of the present proofs.

\bibitem{Sutherland21}
A.~V. Sutherland, \emph{18.785 Number Theory I}, Lecture~28,
Massachusetts Institute of Technology, Fall 2021,
Section~28.4, Theorems~28.6 and~28.9, and Remark~28.12.
\url{https://math.mit.edu/classes/18.785/2021fa/LectureNotes28.pdf}.

\end{thebibliography}

\end{document}
